% concatenated from GPM-est2.tex
\documentclass[a4paper,10pt,twoside]{article}


\pagestyle{plain}

\usepackage{amsmath, amssymb, amsthm, ascmac}
\usepackage{color}
\newcommand{\red}{\color[rgb]{0.8,0,0}}
\newcommand{\blue}{\color[rgb]{0,0,1}}
\usepackage{graphicx}

\usepackage{mathptmx}
%times,fouriernc,fourier,mathptmx,pxfonts,txfonts,euler,serkis
%\usepackage{refcheck}
%\renewcommand{\baselinestretch}{2.0}

\setlength{\oddsidemargin}{5mm}
\setlength{\evensidemargin}{5mm}
\setlength{\textwidth}{150mm}
\setlength{\topmargin}{5mm}
\setlength{\textheight}{210mm}
\setlength{\headsep}{15mm}
\setlength{\headheight}{5mm}


%\theoremstyle{definition}

\newtheorem{thm}{Theorem}[section]
\newtheorem{prop}{Proposition}[section]
\newtheorem{lem}{Lemma}[section]
\newtheorem{cor}{Corollary}[section]
\newtheorem{remark}{Remark}[section]
\newtheorem{example}{Example}[section]
\newtheorem{defn}{Definition}[section]
\newtheorem{asA}{A}
\newtheorem{as}{B}

\numberwithin{equation}{section}
\allowdisplaybreaks

%\renewcommand{\figurename}{Fig.}
%\renewcommand{\tablename}{Tab.}


\def\dis{\displaystyle}
\def\R{\mathbb{R}}
\def\Z{\mathbb{Z}}
\def\C{\mathbb{C}}
\def\N{\mathbb{N}}
\def\H{\mathbb{H}}
\def\Y{\mathbb{Y}}
\def\ve{{\varepsilon}}
\def\e{{\epsilon}}
\def\D{\Delta}
\def\d{\delta}
\def\l{\left}
\def\r{\right}
\def\a{\alpha}
\def\b{\beta}
\def\G{\Gamma}
\def\g{\gamma}
\def\Th{\Theta}
\def\th{\theta}
\def\S{\Sigma}
\def\s{\sigma}
\def\z{\zeta}
\def\k{\kappa}
\def\la{\lambda}
\def\vr{\varrho}
\def\vp{\varphi}
\def\n{\nabla}
\def\p{\partial}
\def\F{\mathcal{F}}
\def\A{\mathcal{A}}
\def\B{\mathcal{B}}
\def\x{\mathcal{X}}
%\def\S{\mathcal{S}}
\def\P{{\mathbb{P}}}
\def\toP{\stackrel{p}{\longrightarrow}}
\def\toD{\stackrel{{\cal L}}{\longrightarrow}}
\def\toSt{\stackrel{{\cal L}_{st}}{\longrightarrow}}
\def\E{\mathbb{E}}
\def\mb#1{\mbox{\boldmath $#1$}}
\def\zero{\mathbf{O}}
\def\I{\mathbf{1}}
\def\wh#1{\widehat{#1}}
\def\ol#1{\overline{#1}}
\def\wt#1{\widetilde{#1}}
\def\Var{\mathrm{Var}}
\def\Cov{\mathrm{Cov}}
\def\df{\mathrm{d}}


\newcommand{\bi}{\begin{itemize}}
\newcommand{\ei}{\end{itemize}}


\title{M-estimation for Gaussian processes with time-inhomogeneous drifts from high-frequency data}
\author{Yasutaka SHIMIZU\\ {\it Department of Applied Mathematics, Waseda University}}
\date{\today}


\begin{document}
\maketitle

\begin{abstract} 
We propose a contrast-based estimation method for Gaussian processes with time-inhomogeneous drifts under high-frequency sampling.
The method constructs a local-Gauss contrast from first-order increments, avoiding inversion of large covariance matrices and allowing computationally tractable inference.
We establish consistency and asymptotic normality under general ergodicity conditions for the driving Gaussian processes.
A key feature of the framework is that the local contrast asymptotically depends only on the local curvature of the covariance kernel around the origin.
Consequently, first-order local increments generally provide only partial information on the covariance structure, and additional kernel parameters are recovered through auxiliary moment-based procedures.
The drift estimator exhibits a nonstandard convergence rate under directly Riemann integrable (DRI) drift structures, reflecting weak information accumulation even when the observation horizon diverges.
The framework is flexible and applicable to a broad class of Gaussian process models, including nonsmooth kernels through mollified approximations.
Several examples and numerical experiments illustrate the theoretical and finite-sample properties of the proposed estimators.

\begin{flushleft}
{ \it Keywords:} Gaussian processes, high-frequency data, time-inhomogeneous drift, contrast-based estimation, method of moments
\vspace{1mm}\\
{\it MSC2010:} {\bf 62M10}; 62F12,  60G15.
\end{flushleft}
\end{abstract} 


\section{Introduction}

Gaussian processes (GPs) provide a flexible probabilistic framework for modeling time-dependent phenomena with uncertainty quantification.
Because of their analytical tractability and nonparametric structure, they are widely used in statistics, machine learning, spatial statistics, and time series analysis. 
In many practical situations, observed data exhibit both deterministic long-term trends and stochastic short-term fluctuations.
This naturally leads to models consisting of a time-inhomogeneous drift combined with a stationary Gaussian component.

In this paper, we consider a stochastic process of the form
\[
X_t = Z_t + \int_0^t \mu_\xi(s)\mathrm{d}s,
\qquad t \ge 0,
\]
where $\mu_\xi$ is a deterministic drift density, $Z=(Z_t)_{t\ge0}$ is a centered stationary Gaussian process
with covariance kernel $K_\sigma$, defined as follows: 

\begin{defn}
A stochastic process $Z = (Z_t)_{t \ge 0}$ is called a \textit{Gaussian process} if and only if any finite-dimensional distribution is multivariate Gaussian: for any $t_1,\dots,t_d > 0$, there exist a mean vector $m \in \mathbb{R}^d$ and a positive definite $d \times d$ matrix $\Sigma$ such that
\[
(Z_{t_1}, Z_{t_2}, \dots, Z_{t_d}) \sim N_d(m, \Sigma),
\]
where $N_d(m, \Sigma)$ denotes the $d$-dimensional normal distribution with mean $m$ and covariance matrix $\Sigma$.
\end{defn}

\begin{defn}
A Gaussian process $Z$ is said to be \textit{stationary} if there exists a function $K:[0,\infty) \to \mathbb{R}$ such that
\[
\E[Z_t Z_s] = K(|t - s|), \quad t, s \ge 0.
\]
That is, the covariance between $Z_t$ and $Z_s$ depends only on their time lag. The function $K$ is called the \textit{kernel function}. Moreover, $Z$ is said to be \textit{centered} if the mean function $m(t) := \E[Z_t]$ is identically zero.
\end{defn}


The process is observed under high-frequency sampling:
\[
t_i = ih_n,
\qquad
h_n \to 0,
\qquad
nh_n \to \infty.
\]
Our objective is to estimate both the drift parameter $\xi$ and the covariance parameter $\sigma$ from discrete observations of $X$.

Parameter estimation for Gaussian processes has been extensively studied.
The most classical approach is maximum likelihood estimation (MLE), which relies on the multivariate Gaussian likelihood; see, e.g., Rasmussen and Williams~\cite{rw06}.
However, naive evaluation of the Gaussian likelihood requires inversion of large covariance matrices, whose computational complexity is typically of order $O(n^3)$.
This becomes problematic for high-frequency or large-scale data.

To overcome this difficulty, various approximate likelihood methods have been proposed.
Composite likelihood methods, developed by Cox and Reid~\cite{cr04}, Davis and Yau~\cite{dy11}, Varin et al.~\cite{vetal11}, and Bachoc and Lagnoux~\cite{bl20}, construct tractable criteria from low-dimensional Gaussian marginals.
In the context of high-frequency Gaussian processes, Bennedsen et al.~\cite{betal24} proposed composite likelihood estimators based on local Gaussian blocks.
Frequency-domain approaches, originating from Whittle~\cite{w51}, have also been extensively developed; see, for example, Takabatake~\cite{t22} and Fukasawa and Takabatake~\cite{ft16} for high-frequency inference for Gaussian models.
Continuous-time likelihood approaches for Gaussian process models have also recently been investigated by Kobayashi et al.~\cite{ketal25}.

The present paper is also based on local increments.
More precisely, we construct a contrast function from first-order increments $\Delta_i^nX=X_{t_i}-X_{t_{i-1}}$ using only their local Gaussian approximation.
The resulting criterion depends only on scalar local variances, and does not require covariance matrix inversion. 

The primary purpose of the proposed local-Gauss contrast is inference for the drift component from local increment information.
The covariance kernel enters mainly through its local behavior around the origin, which determines the asymptotic structure of first-order increments.
This naturally raises the question of how much covariance information can be recovered from local increments alone.

A key observation is that first-order local increments contain only localized information about the covariance structure.
More precisely, the local contrast asymptotically depends only on the local curvature $\p_t^2K_\sigma(0)$ through the expansion
\[
K_\sigma(h_n)-K_\sigma(0)
=
\frac12\p_t^2K_\sigma(0)h_n^2 + o(h_n^2).
\]
Consequently, first-order local increments generally provide only partial information on the covariance structure.
For instance, in Gaussian kernel models, only certain combinations of covariance parameters appear asymptotically in the local contrast.
Thus, the proposed local-Gauss contrast should not be interpreted as a full covariance estimation procedure.
Instead, additional covariance parameters are recovered through auxiliary moment-type estimators designed to capture global covariance information invisible to local increments.

Another distinctive feature of the present framework is the appearance of a nonstandard asymptotic rate.
Under directly Riemann integrable (DRI) drift structures, the estimator of the drift parameter satisfies
\[
h_n^{-1/2}(\widehat{\xi}_n-\xi_0)=O_p(1),
\]
instead of the standard $\sqrt{nh_n}$-type rate commonly encountered in ergodic diffusion inference.
This phenomenon originates from the fact that the proposed local-Gauss contrast is driven primarily by local increment information.
Under DRI-type assumptions, the effective information accumulation for the drift component remains relatively weak even when the observation horizon
\[
T_n := nh_n
\]
diverges.
Consequently, although consistency and asymptotic normality hold theoretically, finite-sample concentration may remain weak, making the convergence behavior difficult to observe numerically for moderately large sample sizes.

Thus, the main focus of this paper is not merely computational simplification, but rather the asymptotic structure generated by local increment inference.
In particular, the paper reveals:
\begin{itemize}
\item
a partial identifiability structure induced by first-order local increments;

\item
a nonstandard convergence rate caused by DRI-type information accumulation;

\item
an asymptotic separation between local covariance curvature and global covariance information.
\end{itemize}

From a technical viewpoint, the paper also develops several asymptotic tools for ergodic Gaussian processes under high-frequency sampling.
In particular, the appendix establishes weak-dependence limit theorems and moment estimates for local increment functionals under general ergodicity conditions.
These results are not restricted to the present estimation problem, and may provide useful probabilistic tools for future studies of high-frequency inference for non-Markovian Gaussian models.

The asymptotic theory developed in this paper also provides a theoretical basis for future hybrid procedures combining local increment inference with richer covariance recovery methods, including composite likelihood and frequency-domain approaches.

The remainder of the paper is organized as follows. 
Section~\ref{sec:models} introduces the Gaussian process model and assumptions.
Section~\ref{sec:local-Gauss} develops the local-Gauss contrast estimator and establishes consistency and asymptotic normality.
Section~\ref{sec:moment} studies moment-type estimators for recovering additional covariance information.
Section~\ref{sec:example} provides examples and numerical experiments.
Proofs are collected in later sections and appendices.


\subsection*{Notation}

\begin{itemize}
\item The random vector $Z$ follows the normal distribution with mean vector $m$ and covariance matrix $\Sigma$, we write $Z\sim {\cal N}(m,\Sigma)$. 
\item The probability density function of ${\cal N}(0,\Sigma)$ is given by $\phi_{\Sigma}$. 
\item For a centered stationary Gaussian process $Z=(Z_t)_{t\ge 0}$ with the kernel function $K$, we write $Z\sim GP(0,K)$.  
\item For a subet $S\subset \R^p$, $\ol{S}$ is the closure of $S$ w.r.t. the Euclidian norm. 
\item For a function $f(x,y): \R^{d}\times \R^{d'}\to \R$ and $x=(x_1,\dots,x_d)^\top$, 
\[
\p_x f := \frac{\p f}{\p x} = \l(\frac{\p f}{\p x_1},\dots,\frac{\p f}{\p x_d}\r)^\top \in \R^{d},\quad 
\p_x^2 f := \p_x \p_x^\top f = 
\begin{pmatrix}
\frac{\p^2 f}{\p x_1\p x_1} &\dots& \frac{\p^2 f}{\p x_1 \p x_d} \\
\vdots & \ddots & \vdots \\
\frac{\p^2 f}{\p x_d\p x_1} &\dots& \frac{\p^2 f}{\p x_d \p x_d}
\end{pmatrix} \in \R^{d\times d}. 
\]
if the partial derivatives exist, where $\top$ stands for the transpose. 

\end{itemize}


\section{Models and Assumptions}\label{sec:models}

Consider a stochastic process driven by a Gaussian process $Z$:
\begin{align}
X_t = Z_t + \int_0^t \mu(s)\,\df s,\quad  X_0 = Z_0\label{model}
\end{align}
where $\mu: [0,\infty) \to \R$ and $Z=(Z_t)_{t\ge 0}\sim GP(0,K)$, a centered stationary Gaussian process with the kernel function $K$. 
The goal of the paper is to estimate the mean density $\mu$ and kernel functions $K$  from discrete samples of $X$ as follows:  
\[
X_{t_0},\ X_{t_1},\  \dots,\  X_{t_{n-1}},\  X_{t_n}, 
\]
where $t_i = i h_n$ with $h_n > 0$ for $i=1,2,\dots,n$. In asymptotic theory, we assume high-frequency sampling over a long time horizon, which is standard in modern applications where sufficiently dense observations are available: 
\begin{align}
h_n\to 0,\quad nh_n\to \infty,\quad \mbox{as $n\to \infty$. } \label{sampling}
\end{align}

For that purpose, we consider parametric families for $\mu$ and $K$ as follows: 
\[
\l\{\mu_\xi:  [0,\infty) \to \R \,|\, \xi \in \ol{\Xi}\r\};\quad \l\{K_\s:[0,\infty) \to (0,\infty)\,|\, \s\in \ol{\Pi}\r\}, 
\]
where  $\Xi\subset \R^p$  and $\Pi\subset \R^q$ are open and convex bounded subset, and set $\Theta := \Xi \times \Pi$. 
We suppose that there is the true values of parameters: 
\[
\th_0=(\xi_0,\s_0) \in \Theta;\quad \mu_{\xi_0} \equiv \mu;\quad K \equiv K_{\s_0}. 
\]

\begin{asA}\label{as:2}
$K(t) \to 0$ as $t \to \infty$.
\end{asA}

\begin{asA}\label{as:3}
 $K \in C^2([0,\infty))$ with $\p_t K(0) = 0$ and $|\p_t^2 K(0)| > 0$, that is,
\[
K(t) = K(0) + \frac{1}{2} \p_t^2 K(0) t^2 + o(t^3), \quad t \to 0.
\]
\end{asA}

\begin{asA}\label{as:4}
$K \in C^2([0,\infty))$ and 
\[
\p_tK(t) \to 0,\quad t\to \infty. 
\]
\end{asA}
Condition A\ref{as:2} is a mixing-type condition. In fact, the process $Z$ is \textit{weakly mixing} under A\ref{as:2}; see Maruyama~\cite{m49}. This ensures the ergodicity of $Z$ in the sense of Corollary~\ref{cor:ergodic}.

This smoothness condition relates to the regularity of sample paths of the underlying Gaussian process. The assumption A\ref{as:3} is just the Taylor expansion and the remainder will be $o(t^3)$ since $K$ is symmetric. However, such a smoothness is not always standard; in fact, certain important examples such as the Ornstein-Uhlenbeck (O-U) process (see Example~\ref{ex:ou-kernel}) do not satisfy it. Nonetheless, consistent estimators of $\theta$ can be constructed by approximating the non-smooth kernel with a smooth mollified version $K_\varepsilon$ such that $K_\varepsilon \to K$ as $\varepsilon \to 0$; see Section~\ref{ex:O-U}.


We shall give some examples on stationary kernels. 
\begin{example} \label{ex:kernels}
Let us give some examples for kernel functions satisfying A\ref{as:2} and A\ref{as:3}. 
\bi
\item {\it Gaussian kernel (Radial Basis Function)}: 
\[
K_\s(t) = \a \exp\l(-\frac{\b}{2} t^2\r),\quad \s=(\a,\b) \in \R^2,
\]
with 
\[
K_\s(t) = \a - \a\b\,t^2 + o(t^3),\quad t\to 0.
\]

\item {\it Mat\'ern kernel}: (this is of $C^2$ as $\nu > 2$) 
\[
K_\s(t) = \a \frac{2^{1-\nu}}{\G(\nu)} \l( \sqrt{2\nu} \b |t| \r)^\nu B_\nu\l( \sqrt{2\nu} \b |t| \r),\quad \s=(\a,\b,\nu)\in \R^3
\]
where $B_\nu$ is the modified Bessel function of 2nd kind. It is known that 
\[
K_\s(t) =\a - \a \b^2\frac{\nu }{2\nu - 1} t^2 + o(t^3),\quad t\to 0. 
\]
\item {\it Rational Quadratic kernel}
\[
K_\s(t) = \a \l(1 + \frac{\b^2 t^2}{2\g} \r)^{-\g},\quad \s=(\a,\b,\g)\in \R_+^3, 
\]
with 
\[
K_\s(t) = \a - \frac{\a \b^2}{2} t^2 + o(t^3),\quad t\to 0, 
\]
In particular, it is called `Cauchy kernel' as $\a=1$. 
\ei
\end{example}

\begin{example} \label{ex:ou-kernel}{Exponential (O-U) kernel} 
The following kernel is also important in applications. 
It is called the {\it exponentail kernel} or the {\it Ornstein-Uhlenbeck kernel}: 
\[
K_\s(t) = \a  \exp\l(-\b|t|\r),\quad \s=(\a,\b) \in \R_+^2,
\]
which is not smooth at $t=0$ since it includes $|t|$ in the exponent. However, we can approximate this kernel with a {\it `mollifier'}  such as
\[
\int_\R \vp(s)\,\df s = 1; \quad \vp_\e(t):=\e^{-1} \vp\l(\e^{-1} t\r) \to \d_0(t)\quad \e\to 0, 
\]
where $\d_0$ is Dirac's delta function. Consider a {\it smoothed kernel}
\[
K_\s^{(\e)}(t) = \int_\R K_\s(t-s)\vp_\e(s)\,\df s \to K_\s(t) \quad \e\to 0. 
\]
Then, it follows for $m\in \N$ that 
\[
\p_t^m K_\s^{(\e)}(t) := \int_\R K_\s(t-s)\p^m \vp_\e(s)\,\df s \to \p_t^mK_\s(0),\quad \e\to 0, 
\]
where the last $ \p_t^mK_\s(0)$ is a `generalized' derivative. 
Therefore, a smooth $K_\s^{(\e)}$ is available as an approximation of non-smooth $K(t)$ in practice. 
For example, using the Laplace mollifier 
\[
\varphi_\e(s) = \frac{1}{2\e} e^{-\frac{|s|}{\e}}, 
\]
we have for each $t > 0$:
\begin{align*}
K_\s^{(\e)}(t) &= \a e^{-\b t} \int_{\R} e^{\b s} \vp_\e(s) \mathbf{1}_{\{s \le t\}}\,\df s + \a e^{\b t} \int_{\R} e^{-\b s} \vp_\e(s) \mathbf{1}_{\{s > t\}}\,\df s\to \frac{\a}{1 + \b\e}  \quad (t\to 0) \\
\p_t K_\s^{(\e)}(t) &= -\a\b e^{-\b t} \int_{\R} e^{\b s} \vp_\e(s) \mathbf{1}_{\{s \le t\}}\,\df s + \a\b e^{\b t} \int_{\R} e^{-\b s} \vp_\e(s) \mathbf{1}_{\{s > t\}}\,\df s, \\
&\to  -\a\b \int_{-\infty}^0 e^{\b s}\vp_\e(s)\,\df s + \a\b \int_0^\infty e^{-\b s}\vp_\e(s)\,\df s =0,  \quad (t\to 0) \\
\p_t^2 K_\s^{(\e)}(t) &= \a\b^2 \l[ e^{-\b t} \int_{\R} e^{\b s} \vp_\e(s) \mathbf{1}_{\{s \le t\}}\,\df s + e^{\b t} \int_{\R} e^{-\b s} \vp_\e(s) \mathbf{1}_{\{s > t\}}\,\df s \r] - 2\a\b \vp_\e(t) \\
&\to \frac{\a\b^2}{1+\b\e} - \frac{\a\b}{\e}. 
\end{align*}

Moreover, as $t \downarrow 0$, using this explicit form we obtain the following expansion:
\begin{align*}
K_\s^{(\e)}(t) &=\frac{\a}{1 + \b\e} -  \l(\frac{\a\b}{\e} - \frac{\a\b^2}{1+\b\e}\r) t^2 + o(t^3), \quad (t\to 0). 
\end{align*}
Hence, the mollified kernel $K_\s^{(\e)}$ is smooth and satisfies the condition A\ref{as:3} for any fixed $\e > 0$. 

That is, instead of modeling the data by an exact Ornstein-Uhlenbeck process, one can model it by a Gaussian process with kernel $K_\sigma^{(\varepsilon)}$ using a small $\varepsilon > 0$. See the example in Subsection~\ref{ex:O-U} for an illustration of this approach.

\end{example}

\begin{remark}
As for the assumption A\ref{as:3}, there exists a Gaussian process with $\p_t^2 K_\s(0) = 0$. For example, 
\[
K_\s(t) = 
\begin{cases}
\exp\l(-\frac{1}{1 - \l(t/\s\r)^2}\r) & \text{if } |t| < \s, \\
0 & \text{otherwise}, 
\end{cases}
\]
is of $C^\infty$-class, and all the derivatives at $t=0$ is identically zero. Since we can also confirm the semi-positive definiteness of this kernel, 
this is a Gaussian kernel.
However, such a Gaussian process with `flat' derivatives at zero is  impractical because the sample paths are too smooth due to the small quadratic variation. From an applied point of view, we are interested in processes with higher volatilities. Hence, A\ref{as:3} is not so strong assumption. 
\end{remark}


In what follows, we shall use the concept of \emph{directly Riemann integrable (DRI)}: A non-negative function $g : \mathbb{R} \to [0, \infty)$ is said to be DRI if its upper and lower Riemann sums \emph{over the whole real line} converge to the same finite limit, as the mesh of the partition vanishes: 
\[
\lim_{h \downarrow 0} \sum_{k \in \mathbb{Z}}  \underline{g}_{k,h} \cdot h=\lim_{h \downarrow 0} \sum_{k \in \mathbb{Z}}  \ol{g}_{k,h}\cdot h =:\int_\R g(s)\,\df s \in (-\infty,\infty), 
\]
where $\ol{g}_{k,h}:=\sup_{z \in [kh, (k+1)h)} g(z)$ and $\underline{g}_{k,h} := \inf_{z \in [kh, (k+1)h)} g(z)$. 
If the function $g$ may also take negative values, it is said to be DRI if both its positive and negative parts $g^+$ and $g^-$ are so. For more details, see, e.g., Asmussen~\cite{a03}, Section V.4; Feller~\cite{f71}, Section~XI.1; Rolski et al.~\cite{retal99}; Caravenna~\cite{c12}; and references therein. 

\begin{remark}\label{rem:dri-cont}
Note that if $g$ is DRI, then $g$ is bounded and continuous a.e. with respect to the Lebesgue measure; see Asmussen~\cite{a03}, Proposition V.4.1. 
\end{remark}

\begin{remark}\label{rem:DRI}
In our sampling scheme \eqref{sampling}, it follows for a DRI function $g:[0,\infty)\to \R$ that 
\begin{align}
\lim_{n\to \infty}  \sum_{i=1}^n g(t_{i-1})h_n = \int_0^\infty g(s)\,\df s, \label{DRI}
\end{align}
since, under $nh_n\to \infty$, 
\[
\lim_{n\to \infty} \l[\sum_{i=1}^\infty g(t_{i-1})h_n  - \sum_{i=1}^n g(t_{i-1})h_n\r] = \lim_{n\to \infty} \int_{nh_n}^\infty g(s)\,\df s = 0. 
\]
Therefore, when we use the convergence \eqref{DRI}, the condition $nh_n\to \infty$ is always required. 

\end{remark}



We further impose the following assumptions on the parametric model, which will be introduced as needed in the discussion below.

\begin{as}\label{as:K2}
For any $\s\in \ol{\Pi}$,  $K_\s \in C^2([0,\infty))$ with $\p_tK_\s(0) = 0$ and $|\p_t^2 K_\s(0)| \in (0,\infty)$, that is, 
\[
\sup_{\s\in \ol{\Pi}} \l| \frac{K_\s(t) -  K_\s(0) - \frac{1}{2}\p_t^2 K_\s(0) t^2}{t^2}\r| \to 0,\quad t\downarrow 0. 
\]
\end{as}

\begin{as}\label{as:mu}
For any $\xi \in \Xi$, $\mu_\xi \in C^1(\R)$ and $\sup_{t>0,\xi\in \ol{\Xi}}|\p_t \mu_\xi(t)| < \infty$. 
\end{as}

\begin{as}\label{as:DRI-mu}
For any $\xi \in \Xi$, $\mu_\xi$ is DRI on $\R$.
\end{as}

\begin{as}\label{as:id2}
$\mu_\xi = \mu_{\xi_0}\ a.e.$ for any $t\in [0,\infty)$ implies that $\xi = \xi_0$. 
\end{as}

\begin{as}\label{as:id}
$\p^2_t K_\s(0) = \p^2_t K_{\s_0}(0)$ implies that $\s = \s_0$. 
\end{as}

\begin{as}\label{as:DRI-mu2}
For any $\xi \in \Xi$, the function $t \mapsto \partial_\xi \mu_\xi(t)$ is DRI on $[0, \infty)$.
\end{as}

\begin{as}\label{as:DRI-mu4}
For any $\xi \in \Xi$, the function $t \mapsto \partial_\xi^2 \mu_\xi(t)$ is bounded.
\end{as}

\begin{as}\label{as:DRI-mu3}
There exists a bounded function $\g \in L^1([0,\infty))$ such that $\sup_{\xi \in \Xi} |\partial_\xi \mu_\xi(t)|\le \g(t)$ for each $t\ge 0$.
\end{as}

\begin{remark}
Assumption B\ref{as:id} is related to the identifiability properties of the local contrast based on first-order increments.
Since the local contrast depends asymptotically only on the local quadratic behavior of the covariance kernel around the origin, different parameter values may yield the same local curvature structure.

For instance, in the Gaussian kernel model
\[
K_\s(t)=\a\exp\l(-\frac{\b}{2}t^2\r),
\]
the quantity $\p_t^2K_\s(0)=-2\a\b$ depends only on the product parameter $\a\b$.
Hence, the local contrast alone cannot asymptotically distinguish $\a$ and $\b$ separately.

Nevertheless, additional global information, such as moment-type statistics, may still allow separate estimation of covariance parameters.
This motivates the auxiliary estimation procedure introduced later in the paper.\end{remark}






\section{M-estimation for sampled Gaussian processes}\label{sec:M-est}

\subsection{Local-Gauss contrast}\label{sec:local-Gauss}

The local-Gauss contrast utilizes only the local increment structure of the Gaussian process.
More precisely, the resulting contrast depends asymptotically only on the local quadratic behavior of the covariance kernel around the origin.
Therefore, the corresponding $M$-estimator mainly captures local curvature information of the covariance structure,
rather than the full global covariance behavior.  
This local nature is computationally advantageous, since it avoids inversion of large covariance matrices,
but it may lead to partial identifiability of covariance parameters.
To recover additional global covariance information, we later introduce auxiliary moment estimators.


We use the notation that, for a process $X = (X_t)_{t\ge 0}$, 
\[
\D_i^n X := X_{t_i} - X_{t_{i-1}}, 
\]
the increment of $X$ on $(t_{i-1},t_i]\ (i=1,2,\dots,n)$. Note that
\[
\D_i^nX - \mu(t_{i-1}^n) h_n \approx  \D_i^n Z = (1 , -1)\begin{pmatrix}Z_{t_i^n} \\ Z_{t_{i-1}^n} \end{pmatrix} \sim {\cal N}(0, 2[K(0) - K(h_n)]). 
\]

Hence, we will use the following {\it local-Gauss} contrast function:
\begin{align}
\ell_n(\xi,\s)  = \frac{1}{n}\sum_{i=1}^n \frac{\l(\D_i^n X - \mu_\xi(t_{i-1}) h_n\r)^2}{2[K_{\s}(0) - K_{\s}(h_n)]} +  \log \l(2h_n^{-2}[K_{\s}(0) - K_{\s}(h_n)]\r). \label{lg-contrast}
\end{align}
Note that $2h_n^{-2}$ in the logarithm of the second term is for an appropriate scaling to obtain a proper limt of $\ell_n(\th)$. 
 
We consider a {\it minimum contrast estimator} ($M$-estimator) defined as follows:
\begin{align*}
\wh{\th}_n := (\wh{\xi}_n,\wh{\s}_n)= \arg\min_{\th \in \ol{\Theta}} \ell_n(\th). 
\end{align*}


\begin{remark}\label{rem:xi}
To estimate $\xi_0$, we can optimize the following simplified estimating function:
\[
\wh{\xi}_n=\arg\min_{\xi\in \Xi} \sum_{i=1}^n \l( \Delta_i^n X - h_n \mu_\xi(t_{i-1}) \r)^2,
\]
which corresponds to the least squares estimation, and it does not require an estimator for $\sigma_0$. 
This often yields an explicit form for $\widehat{\xi}_n$; see, e.g., Example~\ref{sec:drift+gauss}.
\end{remark}


\begin{remark}[Relation to composite likelihood methods]
Our contrast function is constructed from the sequence of increments $\D_i^n X := X_{t_i} - X_{t_{i-1}}$, based on the assumption that these are approximately centered Gaussian with variance $K_\sigma(0) - K_\sigma(h_n)$. This yields a pseudo-likelihood that is computationally efficient, as it avoids inversion of covariance matrices.
This approach shares a structural similarity with the composite likelihood estimation (CLE) method proposed by Bennedsen et al.~\cite{betal24}, which constructs a contrast by aggregating marginal Gaussian likelihoods over $q$-dimensional vectors. In particular, when $q = 2$, CLE uses the full bivariate Gaussian likelihood of $(X_{t_{i-1}}, X_{t_i})^\top$, incorporating the full $2 \times 2$ covariance matrix
\[
\Gamma_\sigma^{(2)} =
\begin{bmatrix}
K_\sigma(0) & K_\sigma(h_n) \\
K_\sigma(h_n) & K_\sigma(0)
\end{bmatrix}.
\]
Thus, while both methods are based on local Gaussian structures, our method relies only on scalar increments and is not a strict special case of CLE.
\end{remark}

\begin{thm}\label{thm:consist}
Suppose the assumptions A\ref{as:2}--A\ref{as:4} and B\ref{as:K2}--B\ref{as:id}. 
Then the $M$-estimator $\wh{\th}_n$ 
is consistent to $\th_0$: 
\[
\wh{\th}_n \toP \th_0,\quad n\to \infty.  
\]
under the sampling \eqref{sampling}. 
\end{thm}


\begin{remark}[Local identifiability and auxiliary estimation]
The local-Gauss contrast exploits only the second-order local behavior of the covariance kernel around the origin.
Indeed, under Assumption~B\ref{as:K2},
\[
2[K_\sigma(0)-K_\sigma(h_n)]
=
-\partial_t^2 K_\sigma(0) h_n^2 + o(h_n^2),
\]
and therefore the asymptotic structure of the contrast depends only on the local curvature quantity $-\partial_t^2 K_\sigma(0)$.

Consequently, the estimator $\widehat{\sigma}_n$ generally identifies only combinations of parameters determining
the local curvature of the covariance kernel. 
For example, in Gaussian kernel models, different parameter values may produce the same value of
$\partial_t^2 K_\sigma(0)$.

To recover the remaining covariance parameters, additional global information is necessary.
This motivates the moment-based estimators introduced in the next section, which utilize global covariance structures beyond local increments.
\end{remark}



\begin{thm}\label{thm:asy-norm}
Suppose the same conditions as in Theorem \ref{thm:consist}. Suppose further B\ref{as:DRI-mu2} and B\ref{as:DRI-mu4}. 
Then, the estimator $\wh{\theta}_n$ is asymptotically normal under \eqref{sampling}:
\[
D_n(\wh{\theta}_n - \theta_0) \toD \mathcal{N}(0, \Sigma(\theta_0)),
\]
where $D_n := \mathrm{diag}\l(h_n^{-1/2} I_p, \sqrt{n} I_q\r)$ and the information-type matrix is
\[
\Sigma(\theta_0) := 
\begin{pmatrix}
\dis -\p_t^2 K_{\s_0}(0) \l[2 \int_0^\infty \{\p_\xi \mu_{\xi_0}(t)\}^{\otimes 2}\, \df t\r]^{-1} & 0 \\
0 & V_2^{-1}(\s_0) V_1(\s_0) V_2^{-1}(\s_0)
\end{pmatrix}.
\]
where 
\[
V_1(\s) = \l(\frac{1}{2}\p_\s \log \l(-\p_t^2 K_{\s}(0)\r)\r)^{\otimes 2},\quad V_2(\s) =  \p_\s^2 \log \l(-\p_t^2 K_{\s}(0)\r). 
\]
\end{thm}

\begin{remark}
The rate of convergence for the estimator $\widehat{\xi}_n$ is not the standard $\sqrt{n h_n}$, but rather the nonstandard rate $h_n^{-1/2}$.
Since $\sqrt{n h_n}/h_n^{-1/2} = \sqrt{n h_n^2}$, this rate is faster under the usual high-frequency condition $n h_n^2 \to 0$.
However, this does not imply that $\xi_0$ is estimable from observations over a bounded time interval under the sole condition $h_n \to 0$.
In fact, the condition $n h_n \to \infty$ is essential, and the nonstandard rate originates from the DRI property, reflecting the accumulation of local information through high-frequency sampling; see Remark~\ref{rem:DRI}.
Under this regime, extending the observation horizon does not increase the available information, which is intrinsic to inference based on first-order increments on bounded intervals.
By contrast, in frameworks where standard Riemann approximation applies, such as
\[
\frac{1}{n} \sum_{i=1}^n \mu_\xi(t_{i-1}) \to \int_0^1 \mu_\xi(s) \,\mathrm{d}s \quad (n \to \infty),
\]
one expects the standard convergence rate $\sqrt{n h_n}$ to be recovered, as is typically the case in ergodic or small-noise models.
\end{remark}


\subsection{Moment estimators}\label{sec:moment}
As discussed in Section~\ref{sec:local-Gauss} the local-Gauss contrast exploits only the local quadratic structure
of the covariance kernel through first-order increments.
Consequently, some covariance parameters may remain unidentified when different kernels share the same local curvature behavior.

To recover additional information on the covariance structure, we now introduce auxiliary moment estimators based on
global marginal statistics of the de-trended process.
Unlike the local-Gauss contrast, these estimators utilize the stationary distribution of the Gaussian component
and can identify parameters that do not appear in the local increment asymptotics.

Suppose that an estimator of $\xi_0$ is given, say $\wh{\xi}_n$; see Remark \ref{rem:xi}, and let 
\[
Y_i^n = X_{t_i} - \int_0^{t_i} \mu_{\wh{\xi}_n}(s)\,\df s,\quad i=1,2,\dots,n. 
\]
Consider the following $\R^q$-valued estimating functions: for $f=(f_1,\dots,f_q):\R\to \R^q$, 
\begin{align*}
\Phi_n(\s) &= \frac{1}{n}\sum_{i=1}^n f(Y_{i-1}^n) - \int_\R f(z)\phi_{K_\s(0)}(z)\,\df z. %\label{M-est-func}
\end{align*}
The $Z$-estimator is given by 
\begin{align}
\Phi_n(\wt{\s}_n) = 0\quad (k=1,2). \label{Z-est}
\end{align}
Since it follows by Lemmas \ref{lem:fY} that, for suitable functions $f$ and $G$, 
\begin{align*}
\Phi_n(\s) \toP \Phi(\s)&:=  \int_{\R^q}f(z)\l[\phi_{K_{\s_0}(0)}(z) - \phi_{K_\s(0)}(z)\r]\,\df z;
\end{align*}
as $n\to \infty$ uniformly in $\s\in \ol{\Pi}$. 
Then, $\wt{\s}_n$ can be consistent to $\s_0$ under suitable reguralities. 



\begin{thm}\label{thm:M-consist}
Let $f : \R \to \R^q$ be a measurable function such that $f \in C^1(\R)$ and there exists $C > 0$ such that
\[
|f(x)| + |\p_x f(x)| \lesssim 1 + |x|^C.
\]
Suppose the assumptions A\ref{as:2}--A\ref{as:4} and B\ref{as:DRI-mu3} hold, and that a consistent estimator $\widehat{\xi}_n \toP \xi_0$ is given. 
Moreover, suppose the following identifiability condition is satisfied:
\begin{align}
\inf_{\sigma \in \overline{\Pi} : |\sigma - \sigma_0| > \varepsilon} |\Phi(\sigma)| > 0 \quad \text{for all } \varepsilon > 0.
\label{Phi-identifiable}
\end{align}
Then the $Z$-estimator $\widehat{\sigma}_n$ defined by \eqref{Z-est} is consistent:
\[
\wt{\s}_n \toP \sigma_0, \quad n \to \infty. 
\]
\end{thm}

\begin{remark}[Two-stage interpretation and Assumption B\ref{as:id}] 
Theorem \ref{thm:M-consist} provides preliminary estimation of a subset of parameters via the moment method, which can be regarded as nuisance parameters in the subsequent analysis.
After this reduction, Assumption B\ref{as:id} is imposed only on the remaining key parameter and characterizes the information content of first-order increments under high-frequency sampling.
Under this reduced parameterization, simultaneous estimation is feasible and leads to the contrast-based M-estimator studied in Section \ref{sec:M-est}.
The apparent one-dimensional restriction induced by Assumption B\ref{as:id} therefore reflects an intrinsic limitation of first order increments rather than a vacuity of the inference procedure.
\end{remark}

\begin{thm}\label{thm:M-asy-norm}
Suppose the same assumptions as in Theorem~\ref{thm:M-consist}, and that the function $f : \mathbb{R} \to \mathbb{R}^q$ is uniformly bounded and of polynomial growth.
Suppose further that the limiting function $\Phi : \mathbb{R}^q \to \mathbb{R}^q$ is continuously differentiable at $\sigma_0 \in \Pi$, 
and that the Jacobian matrix $A := \p_\sigma \Phi(\sigma_0) \in \mathbb{R}^{q \times q}$ is invertible.
Moreover, suppoe that the following limit exists:
\[
\Gamma^2 := \lim_{n \to \infty} n\, \mathrm{Var}\l( \Phi_n(\sigma_0) \r) \in \mathbb{R}^{q \times q}.
\]

Then, the asymptotic normality holds true:
\[
\sqrt{n} \l( \wt{\sigma}_n - \sigma_0 \r) 
\xrightarrow{d} \mathcal{N} \l( \mb{0},\ A^{-1} \Gamma^2 A^{-\top} \r), \quad n \to \infty. 
\]
\end{thm}

\begin{remark}[Alternative moment-based estimators using paired observations]
Beyond the moment function $\Phi_n(\sigma)$ based on single-time statistics, one can construct alternative moment-type estimators using local pairs of de-trended observations $(Y_{i-1}^n, Y_i^n)$. Lemma~\ref{lem:GY} provides a general convergence result for statistics of the form
\begin{align*}
\frac{1}{n h_n^2} \sum_{i=1}^n G(Y^n_{i-1}, Y^n_i)&\toP  \frac{\p^2 K(0)}{2 K(0)}  \int_\R \l[\p_y G(z,z)\,z - \p_y^2 G(z, z)K(0)\r] \, \phi_{K(0)}(z)\,\df z,
\end{align*}
uniformly in $\th \in \ol{\Theta}$ under $h_n\to 0$ as $n\to \infty$, where $Y_i^n= X_{t_i} - \int_0^{t_i} \mu_{\wh{\xi}_n}(s)\,\df s$ and 
$\wh{\xi}_n$ is a consistent estimator for $\xi_0$: $\wh{\xi}_n\toP \xi_0$. 
The function $G$ is a smooth function of polynomial growth with $G(x,x)=0$. For example:
\begin{itemize}
  \item $G(x,y) = (y - x)^2$ yields the second-order increment moment. 
  \item $G(x,y) = (y - x) y^3$ captures nonlinear interactions between local variation and the magnitude. 
\end{itemize}
Such functionals can be used to construct moment equations for estimating kernel parameters. Consistency follows directly from Lemma~\ref{lem:GY}; see Section \ref{sec:drift+gauss}. 
\end{remark}


\begin{remark}[Estimation of the fourth derivative $\p_t^4 K(0)$]\label{rem:K4}
In the moment-based approach discussed so far, only the quantities $K(0)$ and $\p_t^2 K(0)$ can be extracted directly. However, if the kernel $K$ depends on three or more parameters, as in the case of rational quadratic kernels, higher-order information such as $\p_t^4 K(0)$ becomes essential for parameter identification and estimation. 
Under $K\in C^4([0,\infty))$, the following procedure provides a general and practical way to incorporate such information.

Let us define $\d := -\p_t^2 K(0) > 0$, so that
\[
\mathbb{E}[(\D_i^n Z)^2] = 2\{K(0) - K(h_n)\} = \d h_n^2 - \frac{1}{12} \p_t^4 K(0) h_n^4 + o(h_n^4), \quad h_n \to 0.
\]
Noting that $\D_i^n Z \sim \mathcal{N}(0, \mathbb{E}[(\D_i^n Z)^2])$, we obtain by Gaussianity
\[
\mathbb{E}[(\D_i^n Z)^4] = 3 \left( \mathbb{E}[(\D_i^n Z)^2] \right)^2 
= 3 \d^2 h_n^4 - \frac{1}{2} \d \p_t^4 K(0) h_n^6 + o(h_n^6).
\]
Hence, the fourth derivative $\p_t^4 K(0)$ is identified by
\[
\p_t^4 K(0) = \frac{2}{\d h_n^2} \left\{3 \d^2 - \mathbb{E}\l[\l(\frac{\D_i^n Z}{h_n}\r)^4\r]  \right\} + o(1).
\]
Then, by Lemma~\ref{lem:f-DX}, the quantity $\d = -\p_t^2 K(0)$ is consistently estimated by
\[
\wh{\d}_n := \frac{1}{n h_n^2} \sum_{i=1}^n \left( \D_i^n X - \mu_{\widehat{\xi}_n}(t_{i-1}) h_n \right)^2.
\]
Using this, we define the estimator of $\p_t^4 K(0)$ as
\begin{align}
\widehat{\p_t^4 K(0)} := \frac{1}{h_n^2} \left\{3 \wh{\d}_n - \frac{1}{n h_n^4 \wh{\d}_n} \sum_{i=1}^n \left( \D_i^n X - \mu_{\widehat{\xi}_n}(t_{i-1}) h_n\right)^4  \right\}.\label{est:K4}
\end{align}
Then, this estimator is consistent as $n \to \infty$, $h_n \to 0$, and $n h_n \to \infty$.
\end{remark}

\section{Examples and simulations} \label{sec:example}


\subsection{Drifted Gaussian processes with Gaussian Kernels}\label{sec:drift+gauss}
Consider a model \eqref{model} with 
\[
\mu_\xi(s) = \xi w(s),\quad K_\s(t) = \a \exp\l(-\frac{\b}{2} t^2\r),
\]
where $w:[0,\infty) \to \R$, is a known function, directly Riemann integrable function, and $\xi \in \R$, $\g:=\a\b\in \R_+^2$. 
Note that, in this model, we can not identify $\a_0$ and $\b_0$ separately, but only $\g_0 := \a_0\b_0$ because $\p^2_t K_\s(0) = - \a\b$. 
Then, our local-Gauss contrast function is given by 
\begin{align*}
\ell_n(\xi,\s) &= \frac{1}{n}\sum_{i=1}^n \frac{\l(\D_i^n X - \xi w_{i-1} h_n \r)^2}{2\a [1 - e^{-\b h_n^2/2}]} + \log \l(2 h_n^{-2}\a[1 - e^{-\b h_n^2/2}]\r) \\
&= \frac{1}{nh_n^2 }\sum_{i=1}^n \frac{\l(\D_i^n X - \xi w_{i-1} h_n \r)^2}{\g} + \log \g + o_p(1), 
\end{align*}
where $w_{i-1} = w(t_{i-1})$. 
Hence, we obtain the following $M$-estimator by solving the estimating equation  $\n \ell_n(\th) = 0$: 
\begin{align}
\wh{\xi}_n = \frac{\sum_{i=1}^n w_{i-1}\D_i^nX }{h_n \sum_{i=1}^n w_{i-1}^2},\quad \wh{\g}_n = \frac{1}{nh_n^2}\sum_{i=1}^n \l(\D_i^n X - \wh{\xi}_n w_{i-1} h_n \r)^2, \label{M-est-Gauss1}
\end{align}
which are asymptotically normal estimators for $\xi_0$ and $\g_0$, respectively, with convergence rates
$h_n^{-1/2}$ for $\widehat{\xi}_n$ and $\sqrt{n}$ for $\widehat{g}_n$.

For separate estimation of $\a$ and $\b$, we shall consider the method of moment. 
For example, we can use Lemmas \ref{lem:fY} and \ref{lem:GY} with $f(x,\th)=x^2$ and $G(x,y,\th) = (y-x)^2$, respectively: it follows for 
\[
Y_i^n := X_{t_i} -  \wh{\xi}_n\int_0^{t_i} w(s)\,\df s;\quad \D_i^n Y := Y_i^n - Y_{i-1}^n, 
\]
that 
\begin{align*}
\frac{1}{n}\sum_{i=1}^n (Y_{i-1}^n)^2 \toP \a;\quad 
\frac{1}{nh_n^2}\sum_{i=1}^n (\D_i^n Y)^2 \toP \a\b;\quad 
%\frac{1}{nh_n^2}\sum_{i=1}^n   \D_i^n Y\cdot (Y_{i-1}^n)^2 \toP 3\a^2\b,
\end{align*}
as $n\to \infty$. For example, using the first convergence, we can estimate $\a$ and $\b$ separately by, for example, 
\begin{align}
\wh{\a}_n = \frac{1}{n}\sum_{i=1}^n (Y_{i-1}^n)^2,\quad 
\wh{\b}_n = \frac{\sum_{i=1}^n (\D_i^n Y)^2}{h_n^2\sum_{i=1}^n (Y_{i-1}^n)^2}\ \l(= \frac{\wh{\g}_n}{\wh{\a}_n}\r),  \label{M-est-Gauss2}
\end{align}
both of which are asymptotically normal with the rate $\sqrt{n}$. 

\begin{remark}
In this model, the local-Gauss contrast function depends only on the product $g = \alpha\beta$ through the expansion
\[
K_\sigma(0) - K_\sigma(h_n) = \alpha \l(1 - e^{-\beta h_n^2/2}\r) = \tfrac{1}{2} \alpha\beta h_n^2 + o(h_n^3),
\]
as $h_n \to 0$. Therefore, $\alpha$ and $\beta$ are not separately identifiable from the contrast function alone. This issue is resolved by the method of moments, which utilizes higher-order statistics of the de-trended process $Y_i^n$. 
Alternatively, composite likelihood methods such as Bennedsen et al.~\cite{betal24} use multivariate Gaussian densities over $q$-dimensional blocks (e.g., $q=3$), incorporating multiple lagged covariances like $K(h_n)$ and $K(2h_n)$. This richer information structure enables the separate identification of $\alpha$ and $\beta$ through nonlinear relationships among the covariances.
\end{remark}

\subsection{Ornstein-Uhlenbeck processes}\label{ex:O-U}

Consider the stationary O-U process given by$X_t = Z_t$ with Exponential (Ornstein-Uhlenbeck) kernel 
\[
K_\sigma(t) = \alpha e^{-\beta |t|}, 
\]
and the target parameter is $\theta:= \s = (\alpha, \beta)$. 

Although the Ornstein-Uhlenbeck kernel does not satisfy the smoothness assumption A\ref{as:3},
this example illustrates that our local-Gauss methodology can still be applied through mollified approximations.
Moreover, this example clarifies the connection between the present Gaussian process approach and classical high-frequency inference for diffusion processes.

\subsubsection{Classical diffusion inference}
As is well known, the process $X$ satisfies the following stochastic differential equation: 
\begin{align*}
 \df X_t = -\b X_t\,\df t + \sqrt{2\a\b}\df W_t,  \quad X_0=Z_0,  %\label{ou-sde}
\end{align*}
for a Wiener process $W$. In the context of inference for SDEs, under the asymptotic regime $h_n \to 0$, $n h_n \to \infty$, and $n h_n^3 \to 0$, local-Gaussian contrast by Kessler \cite{k97}: 
\[
l_n(\alpha, \beta) = - \frac{1}{n} \sum_{i=1}^n \l\{
\frac{(\Delta_i^n X + \beta h_n  X_{t_{i-1}})^2}{4\alpha \beta h_n} + \log(4\pi \alpha \beta h_n)
\r\}, 
\]
gives an asymptotically efficient estimator for $\alpha$ and $\beta$. 
If $\alpha$ is known, we can use the following asymptotically equivalent contrast for $\beta$: 
\[
l_n(\alpha, \beta) = - \frac{1}{n} \sum_{i=1}^n \l\{
\frac{(\Delta_i^n X )^2}{4\alpha \beta h_n} + \log(4\pi \alpha \beta h_n) 
\r\}. 
\]
Now, we suppose that $\a$ is known. Then an asymptotic efficient estimator of $\b$ given by the maximizer of the last contrast function: 
if $\a$ is known, then 
\begin{align}
\widetilde{\b}_n &= \frac{1}{4 \a nh_n} \sum_{i=1}^n (\Delta_i^n X)^2,\quad n\to \infty. 
\label{est-sde}
\end{align}
and it holds that
\[
\sqrt{n} (\widetilde{\b}_n - \b) \toD \mathcal{N}(0, 2\b^{2}), 
\]
which is an benchmark of the estimator. 


\subsubsection{Mollified local-Gauss inference}

We shall reconsider this estimation from the view point of Gaussian processes. 
First, we can use Lemma \ref{lem:ergod} with $f(z)=z^2$ to estimate $\a$: 
\begin{align*}
\wh{\a}_n = \frac{1}{n}\sum_{i=1}^n (X_{t_{i-1}})^2\toP \a, %\label{alpha-e}
\end{align*}
which is also asymptotically normal by Theorem~\ref{thm:M-asy-norm}: 
\[
\sqrt{n}(\wh{\a}_n- \a) \toD {\cal N}\l(0,2\a^2\r),\quad n\to \infty. 
\]


Because the O-U kernel is not smooth as in A\ref{as:3}, we will use a mollifier $\varphi_\e$ with $\int \varphi_\e(s)\,\mathrm{d}s = 1$, and approximate $K$ by the smoothed kernel as in Example \ref{ex:ou-kernel}: 
\[
K_\sigma^{(\e)}(t) := \int_{\mathbb{R}} K_\sigma(t - s) \varphi_\e(s) \, \mathrm{d}s.
\]
That is, instead of modeling the data by an exact Ornstein-Uhlenbeck process, one can model it by a Gaussian process with kernel $K_\sigma^{(\varepsilon)}$ using a small $\varepsilon > 0$. 

The contrast function is
\[
\ell_n^{(\e)}(\alpha, \beta) = \frac{S_n}{2 V^{(\e)}(\alpha, \beta)} + \log\l(2 h_n^{-2} V^{(\e)}(\alpha, \beta)\r).
\]
where $S_n := \frac{1}{n} \sum_{i=1}^n (\Delta_i^n X)^2$ and $V^{(\e)}(\sigma) := K_\sigma^{(\e)}(0) - K_\sigma^{(\e)}(h_n)$. 
The score function is given by 
\begin{align*}
\frac{\partial \ell_n^{(\e)}}{\partial \alpha} &= \l( 1 -\frac{S_n}{2 V^{(\e)}(\a,\b)} \r)\cdot \frac{\partial}{\partial \a}\log V^{(\e)}(\a,\b), \\
\frac{\partial \ell_n^{(\e)}}{\partial \beta} &=\l( 1 - \frac{S_n}{2 V^{(\e)}(\a,\b)} \r) \cdot \frac{\partial}{\partial \beta}\log V^{(\e)}(\a,\b).
\end{align*}
Therefore an M-estimator is given by solving the equation 
\[
V^{(\e)}(\a,\b) = \frac{1}{2}S_n.
\]

To obtain an explicit estimator we shall use the {\it Laplace mollifier} $\varphi_\e(s) = \frac{1}{2\e} e^{-\frac{|s|}{\e}}$. 
Then, thanks to Example \ref{ex:ou-kernel}, 
\begin{align*}
K_\s^{(\e)}(t) &=\frac{\a}{1 + \b\e} + \l(\frac{\a\b^2}{1+\b\e} - \frac{\a^2\b}{\e}\r) t^2 + o(t^3), \quad (t\to 0) 
\end{align*}
and, as $\ve^{-1} > \b$, 
\begin{align*}
V^{(\e)}(\a,\b) &= \l( \frac{\alpha^2 \beta}{\e} - \frac{\alpha \beta^2}{1 + \beta \e} \r) h_n^2 + o(h_n^3) = \a\b \frac{h_n^2}{\e} + o\l(\frac{h_n^3}{\e}\r). 
\end{align*}
For example, if $\b$ is known then $\a$ is identifiable (satisfying B\ref{as:id}), and we have 
\begin{align*}
\widehat{\b}_n^{(\e)} &= \dfrac{1}{2} S_n\l[ \wh{\a}_n \frac{h^2}{\e} + o\l(\frac{h^3}{\e}\r) \r]^{-1} = \frac{\e}{2\wh{\a}_n nh_n^2} \sum_{i=1}^n (\D_i^n X)^2. %\label{beta-e}
\end{align*}
Hence, taking 
\begin{align*}
\e = h_n/2\ (\to 0), %\label{epsilon-choice}
\end{align*}
our estimator `formally' recovers the asymptotically efficient variance appearing in \eqref{est-sde}. 
Moreover, Theorem \ref{thm:asy-norm} leads that: for any fixed $\e>0$, 
\begin{align*}
\sqrt{n}(\widehat{\b}_n^{(\e)} - \b) \toD \mathcal{N}\l(0, 2\b^{2} + r_\e\r),\quad n\to \infty,  
\end{align*}
and $r_\e = O(\e)$ as $\e\to 0$. Therefore, our mollified estimator $\widehat{\b}_n^{(\e)}$ is matches the efficient benchmark up to an error of order $O(\varepsilon)$. 

\subsection{The Rational Quadratic kernel}
Consider an example that $Z$ has the Rational Quadratic (RQ) kernel given in Example \ref{ex:kernels}:
\[
K_\s(t)  = \a \l(1 + \frac{\b^2 }{2\g}t^2 \r)^{-\g},\quad \s=(\a,\b,\g)\in \R_+^3. 
\]
This kernel arises as a scale mixture of squared exponential kernels and is widely used in Gaussian process modeling for its flexibility.

The spectral density associated with this kernel has a closed-form expression:
\[
f(\omega) = \a \cdot \frac{\sqrt{2\pi} \, \Gamma(\gamma + \frac{1}{2})}{\Gamma(\gamma)} \cdot \frac{1}{(\b^2 \gamma)^{1/2}} \l(1 + \frac{2\pi^2 \omega^2}{\b^2 \gamma} \r)^{-(\gamma + \frac{1}{2})},
\]
where $\Gamma(\cdot)$ denotes the gamma function. 
While analytically available, this density is nonlinear in all parameters and requires nontrivial numerical treatment for Whittle-type likelihood inference.

In contrast, our method only relies on the second-order behavior of the kernel at the origin. Specifically, a simple Taylor expansion yields:
\begin{align*}
K_\sigma(t) = \a \l\{1 - \frac{\b^2}{2}t^2  + \frac{(1+\gamma) \b^4 }{8 \gamma} t^4 + o(t^4)\r\}, \quad t \to 0. %\label{K-RQ}
\end{align*}
so that
\[
-\p^2_t K_\s(0) =\a \b^2 =:\d.
\]
This quantity enters directly into the contrast function and can be computed in closed form regardless of $\alpha$, allowing for explicit and numerically stable estimation. Therefore, even for kernels with analytically known but numerically complex spectral densities, our contrast-based method offers practical advantages in terms of implementation and stability.

These estimators can be constructed explicitly as follows: To estimate $\xi$ and $\d:=\a\b^2$, the contrast function is given by
\begin{align*}
\ell_n(\xi,\d) = \frac{1}{n}\sum_{i=1}^n \frac{\l(\D_i^n X - \mu_\xi(t_{i-1})h_n\r)^2}{2\d h_n^2}  + \log (2 \d), 
\end{align*}
and, by minimizing this contrast function, we obtain
\[
\wh{\xi}_n = \arg\min_{\xi\in \ol{\Xi}} \sum_{i=1}^n \l(\D_i^n X - \mu_\xi(t_{i-1})h_n\r)^2;\quad 
\wh{\d}_n = \frac{1}{2nh_n^2} \sum_{i=1}^n \l(\D_i^n X - \mu_{\wh{\xi}_n}(t_{i-1})h_n\r)^2. 
\]
Hence we also obtain that, for $Y_i^n := X_{t_i} - \int_0^{t_i} \mu_{\wh{\xi}_n}(s)\,\df s$, 
\[
\wh{\a}_n:= \frac{1}{n}\sum_{i=1}^n (Y_{i-1}^n)^2;\quad \wh{\b}_n = \sqrt{ \frac{\wh{\d}_n}{\wh{\a}_n} }. 
\]
To estimate $\g$, we need the information about $\p_t^4 K_\s$, and we may use \eqref{est:K4} in Remark \ref{rem:K4}: 
\[
\widehat{\p_t^4 K(0)} := \frac{1}{h_n^2} \left\{3 \wh{\d}_n^2 - \frac{1}{nh_n^4 \wh{\d}_n} \sum_{i=1}^n \left( \D_i^n X - \mu_{\wh{\xi}_n}(t_{i-1}) h_n \right)^4 \right\}. 
\]
Noticing that 
\[
\p_t^4 K_\sigma(0) = \frac{3 \a \b^4 (1 + \g)}{\g},
\]
we have the following consistent estimator 
\[
\wh{\g}_n = \frac{3 \wh{\a}_n \wh{\b}_n^4}{\wh{\p_t^4 K(0)} - 3 \wh{\a}_n \wh{\b}_n^4}.
\]

This example illustrates that higher-order local information of the covariance kernel can be systematically incorporated
through higher-order increment moments. 
While the local-Gauss contrast itself only captures the second-order curvature parameter $-\partial_t^2 K(0)$,
additional kernel parameters can be recovered through higher-order derivatives such as $\partial_t^4 K(0)$.
This demonstrates the flexibility of the proposed framework for multi-parameter Gaussian process models.


\subsection{Numerical experiments}\label{sec:num}

The purpose of the following experiments is to illustrate three different aspects
of the proposed inference procedure: (i) asymptotic normality under favorable sampling schemes,
(ii) propagation of estimation errors in joint estimation, and (iii) deterioration of finite-sample normality
when the effective observation horizon $T_n := nh_n$ grows slowly. 
 
\subsubsection{Drifted Gaussian processes with Gaussian Kernels}
Let us consider Example \ref{sec:drift+gauss}:
\[
\mu_\xi(s) = \xi e^{-s},\quad K_\s(t) = \a \exp\l(-\frac{\b}{2} t^2\r),
\]
with the true values of the parameter 
\[
(\xi_0,\a_0,\b_0) = (2.0, 1.0, 1.0). 
\]
We compute the estimators given in \eqref{M-est-Gauss1} and \eqref{M-est-Gauss2}: 
\[
\wh{\xi}_n = \frac{\sum_{i=1}^n e^{-{t_{i-1}}} \D_i^nX }{h_n \sum_{i=1}^n e^{-2{t_{i-1}}}}, \quad 
\wh{\a}_n = \frac{1}{n}\sum_{i=1}^n\l(X_{t_{i-1}} -  \wh{\xi}_n(1 - e^{-{t_{i-1}}})\r)^2,\quad 
\wh{\b}_n =  \frac{\wh{\g}_n}{\wh{\a}_n}, 
\]
where
\[
\wh{\g}_n = \frac{1}{nh_n^2}\sum_{i=1}^n \l(\D_i^n X - \wh{\xi}_n e^{-{t_{i-1}}} h_n \r)^2. 
\]
We shall try the following two cases: 
\bi
\item[(I)] $h_n = n^{-0.4}$, where $T_n :=nh_n = n^{0.6}\to \infty$, and  in estimating $\widehat{\alpha}_n$ or $\widehat{\beta}_n$, the other parameter and $\xi$ were assumed to be known and set to their true values. 

\item[(II)] The same setting as in  Case (I), and all the parameters are estimated jointly (we will use $\wh{\xi}_n$ in estimating $\wh{\a}_n$ and $\wh{\b}_n$). 

\item[(III)] $h_n = n^{-0.8}$, where $T_n:= n^{0.2}\to \infty$, the terminal is smaller than that of (I) and (II). Morever,  in estimating $\widehat{\alpha}_n$ or $\widehat{\beta}_n$, the other parameter and $\xi$ were assumed to be known and set to their true values. 
\ei
For each $n=100,\ 1000,\ 3000$, the experiments are itterated $500$ times, and we shall show the mean and standard deviation (s.d.) for each $\wh{\a}_n,\ \wh{\b}_n$ and $\wh{\xi}_n$ in Tables \ref{tab:gauss1} and  \ref{tab:gauss2}, and normal QQ-plots for each estimators in Figures \ref{fig:QQplot_Gauss1} and \ref{fig:QQplot_Gauss2}, respectively. 

We would like to compare (I) vs. (II), and (I) vs. (III). 

%(I) %%%%%%
\begin{table}[htbp]
\centering
{\bf The result of Case (I)} \\
\begin{tabular}{cccc}
$n$ & $\wh{\xi}_n$ & $\wh{\a}_n$ & $\wh{\b}_n$ \\
\hline
500 & 1.9057 & 1.0294 & 1.0210 \\
    & (1.1316) & (0.3167) & (0.2569) \\
1000 & 1.9059 & 1.0093 & 1.0020 \\
    & (1.1489) & (0.2405) & (0.2029) \\
3000 & 1.9279 & 0.9974 & 1.0106 \\
    & (1.2021) & (0.1620) & (0.1459) \\
\hline
 True & 2.0 & 1.0 & 1.0\\
\hline
\end{tabular}
\caption{Case (I): Means and standard deviations (in parentheses) of the estimators $\widehat{\xi}_n$, $\widehat{\alpha}_n$, and $\widehat{\beta}_n$ over 500 replications, with $h_n = n^{-0.4}$ and other parameters fixed at their true values. The results illustrate good finite-sample accuracy and agreement with the asymptotic normality predicted by theory.}
\label{tab:gauss1}
\end{table}

\begin{figure}[htbp]
  \centering
  \includegraphics[width=1.0\linewidth]{QQplot_Gauss.pdf}
  \caption{Normal QQ plots for Case (I): Scaled estimators $h_n^{-1/2}(\widehat{\xi}_n - \xi_0)$ and $n^{1/2}(\widehat{\alpha}_n - \alpha_0)$, $n^{1/2}(\widehat{\beta}_n - \beta_0)$ over 500 replications. The plots show good agreement with the theoretical normal distribution.}
  \label{fig:QQplot_Gauss1}
\end{figure}
 

%(II) %%%%%%
\begin{table}[htbp]
\centering
{\bf The result of Case (II)} \\
\begin{tabular}{cccc}
$n$ & $\wh{\xi}_n$ & $\wh{\a}_n$ & $\wh{\b}_n$ \\
\hline
500 & 1.8199 & 2.2634 & 0.6274 \\
    & (1.1414) & (1.6769) & (0.3382) \\
1000 & 1.9404 & 2.1516 & 0.6302 \\
    & (1.1125) & (1.5830) & (0.3013) \\
3000 & 1.9476 & 2.1510 & 0.6396 \\
    & (1.0970) & (1.6094) & (0.3028) \\
\hline
 True & 2.0 & 1.0 & 1.0\\
\hline
\end{tabular}
\caption{Case (II): Means and standard deviations (in parentheses) of the estimators $\widehat{\xi}_n$, $\widehat{\alpha}_n$, and $\widehat{\beta}_n$ over 500 replications, with $h_n = n^{-0.4}$ and all parameters estimated jointly. The results show noticeable upward bias in $\widehat{\alpha}_n$ and downward bias in $\widehat{\beta}_n$ due to error propagation from $\widehat{\xi}_n$. }
\label{tab:gauss2}
\end{table}

\begin{figure}[htbp]
  \centering
  \includegraphics[width=1.0\linewidth]{QQplot_Gauss2.pdf}
  \caption{Normal QQ plots for Case (II): Scaled estimators from joint estimation with $h_n = n^{-0.4}$. Upward bias in $\widehat{\alpha}_n$ and downward bias in $\widehat{\beta}_n$ are accompanied by departures from normality, especially for $\widehat{\alpha}_n$ although $\wh{\xi}_n$ still seems to be asymptotically normal.}
  \label{fig:QQplot_Gauss2}
\end{figure}



%(III) %%%%%%
\begin{table}[htbp]
\centering
{\bf The result of Case (III)} \\
\begin{tabular}{cccc}
$n$ & $\wh{\xi}_n$ & $\wh{\a}_n$ & $\wh{\b}_n$ \\
\hline
500 & 1.9828 & 1.0236 & 1.0645 \\
    & (1.1970) & (0.8996) & (0.8531) \\
1000 & 2.0446 & 0.9935 & 1.0318 \\
    & (1.1977) & (0.8225) & (0.7393) \\
3000 & 2.0449 & 1.0202 & 1.0205 \\
    & (1.1730) & (0.7799) & (0.6870) \\
    \hline
 True & 2.0 & 1.0 & 1.0\\
\hline
\end{tabular}
\caption{Case (III):  Means and standard deviations (in parentheses) of the estimators $\widehat{\alpha}_n$ and $\widehat{\beta}_n$ over 500 replications, with $\xi$ fixed at its true value, $h_n = n^{-0.8}$, and other parameters known. Consistency is improved compared to Case (II), but the slow growth of $T_n$ leads to noticeable deviations from normality in finite samples; see Figure \ref{fig:QQplot_Gauss3}, below.}
\label{tab:gauss3}
\end{table}

\begin{figure}[htbp]
  \centering
  \includegraphics[width=1.0\linewidth]{QQplot_Gauss3.pdf}
  \caption{Normal QQ plots for Case (III): Scaled estimators with $\xi$ fixed at its true value and $h_n = n^{-0.8}$. While bias is minimal, the slow growth of $T_n:=nh_n=n^{0.2}$ results in clear deviations from normality compared to Case (I), where $T_n=n^{0.6}$, reflecting insufficient mixing in finite samples.}
  \label{fig:QQplot_Gauss3}
\end{figure}


\subsubsection{Discussion of numerical experiments}

The numerical experiments illustrate several characteristic features of the proposed inference procedure.
In Case (I), where nuisance parameters were fixed at their true values, the estimators exhibited good finite-sample accuracy and agreement with the asymptotic normality predicted by theory.
In particular, the QQ plots indicate that the scaled estimators are already close to Gaussian distributions even for moderate sample sizes.

In contrast, Case (II) showed substantial finite-sample instability due to propagation of estimation errors from $\widehat{\xi}_n$ to $\widehat{\alpha}_n$ and $\widehat{\beta}_n$.
In particular, the estimation error of $\widehat{\xi}_n$ enters the residual-based estimator $\widehat{\alpha}_n$ through squared terms, producing upward bias in $\widehat{\alpha}_n$ and corresponding downward bias in $\widehat{\beta}_n$.
Although this effect decreases asymptotically, it remains non-negligible in finite samples.

Case (III) demonstrated a different phenomenon.
Although the bias was substantially reduced by fixing $\xi$ at its true value, the slow growth of the observation horizon $T_n := nh_n$ led to noticeable deviations from normality in finite samples.
Indeed, under the choice $h_n = n^{-0.8}$, we have $T_n = n^{0.2}$, so the effective ergodic averaging remains weak even for large $n$.
This behavior is fundamentally different from standard ergodic diffusion inference, where the effective information typically grows proportionally to $T_n$.

These results indicate that, under DRI-type drift structures, the finite-sample performance of the estimators depends not only on the formal convergence rates in terms of $h_n$ and $n$, but also on the growth rate of $T_n$.
In particular, weak long-horizon averaging may substantially slow finite-sample concentration despite consistency and asymptotic normality.





\section{Concluding remarks}\label{sec:remark}

In this paper, we proposed a contrast-based inference  for Gaussian processes with time-inhomogeneous drifts observed under high-frequency sampling.
The main idea is to construct a local-Gauss contrast using only adjacent increments and scalar variance structures, thereby avoiding inversion of large covariance matrices.
This leads to estimators that are computationally simple, and theoretically tractable under general weak-dependence and ergodicity conditions.

A central feature of the proposed framework is the separation between local and global information of the covariance structure.
The local-Gauss contrast primarily extracts local curvature information through quantities such as $-\p_t^2 K(0)$.
When this local information alone is insufficient for parameter identification, additional kernel parameters can be recovered through auxiliary moment estimators based on marginal distributions or higher-order increment moments.
In particular, higher-order derivatives such as $\p_t^4 K(0)$ can be estimated consistently from fourth-order residual statistics, allowing inference for multi-parameter kernels including rational quadratic kernels.

The examples and numerical experiments illustrate several characteristic aspects of the methodology.
For smooth kernels such as the Gaussian kernel, the local-Gauss contrast yields explicit estimators with asymptotic normality and favorable finite-sample behavior.

A further characteristic feature of the proposed framework is the appearance of nonstandard convergence rates for the drift parameter under DRI-type drift structures. 
In particular, the estimator of the drift parameter converges at the rate $h_n^{-1/2}$,
which differs fundamentally from the standard $\sqrt{nh_n}$-type behavior typically observed in ergodic diffusion inference.

This phenomenon originates from the fact that the proposed local-Gauss contrast is driven primarily by local increment information. 
Under DRI-type assumptions, the effective information accumulation for the drift component remains relatively weak even when the observation horizon $T_n := nh_n$ diverges.
Consequently, although consistency and asymptotic normality hold theoretically,
finite-sample concentration may remain weak,
making the convergence behavior difficult to observe numerically for moderately large sample sizes.

Another important direction is to relax the DRI-type assumptions on the drift functions.
In the present paper, the DRI condition provides a convenient framework in which the local-Gauss contrast yields consistency and asymptotic normality.
However, it also implies weak information accumulation for the drift component, leading to the nonstandard rate $h_n^{-1/2}$ and slow finite-sample concentration.
For drift functions with non-integrable or polynomial growth structures, such as $\mu_\xi(t)=t^\xi$, the information may accumulate more strongly as $T_n \to \infty$.
In such cases, faster convergence rates than those obtained under the DRI framework may be expected.
A systematic theory for this broader class of time-inhomogeneous drifts will be studied elsewhere.



The simulations also reveal that, under DRI-type drift structures, finite-sample concentration may remain weak even for moderately large sample sizes.
This phenomenon is fundamentally different from standard ergodic diffusion inference, since the information accumulation for the drift parameter is not governed solely by the observation horizon $T_n := nh_n$.
The numerical experiments further demonstrate that joint estimation may induce substantial finite-sample bias propagation from the drift estimator to covariance parameter estimators.

Another important aspect of the present framework is its flexibility with respect to kernel regularity.
Even when the covariance kernel is not smooth at the origin, as in the Ornstein-Uhlenbeck process, mollified approximations allow the local-Gauss methodology to remain applicable.
Moreover, the resulting estimators asymptotically recover the efficient benchmark obtained from classical diffusion-based high-frequency inference.

Compared with frequency-domain approaches such as Whittle-type likelihood methods, the proposed framework works directly in the time domain and naturally accommodates time-inhomogeneous drift structures.
In particular, the method does not require explicit spectral representations, frequency-domain optimization, or preliminary de-trending procedures.
This makes the framework particularly attractive for applications involving complex deterministic trends or covariance kernels without tractable spectral forms.

Several directions remain for future research.
One important problem is to establish a general asymptotic theory for shrinking mollification schemes in nonsmooth kernel models.
Another natural extension is to investigate hybrid procedures combining local-Gauss contrasts with frequency-domain methods after drift removal.
It would also be interesting to study bias-reduction techniques and orthogonalized estimating functions that stabilize joint estimation under DRI-type drift structures.
Finally, extending the framework to broader classes of weakly nonstationary or long-memory Gaussian processes remains an important open problem.


\begin{remark}
The direct Riemann integrability assumption imposed on the drift component is mainly used to guarantee the existence of a deterministic quadratic separation limit of the form
\[
\frac1n
\sum_{i=1}^n
\{\mu_\xi(t_{i-1})-\mu_{\xi_0}(t_{i-1})\}^2
\to
M(\xi,\xi_0).
\]

However, the consistency argument itself does not fundamentally rely on direct Riemann integrability.
More generally, it is sufficient to assume the existence of a positive sequence $a_n \to \infty$ such that
\[
\frac1{a_n}
\sum_{i=1}^n
\{\mu_\xi(t_{i-1})-\mu_{\xi_0}(t_{i-1})\}^2
\to
M(\xi,\xi_0),
\]
uniformly in $\xi$, where
\[
M(\xi,\xi_0)=0
\quad\Longleftrightarrow\quad
\xi=\xi_0.
\]

Under suitable differentiability conditions, the corresponding estimator is expected to satisfy
\[
\widehat{\xi}_n-\xi_0
=
O_p\left\{
\left(
\frac{a_n}{T_n}
\right)^{-1/2}
\right\},
\]
where $T_n=n h_n$.
In the present DRI setting, one typically has $a_n=n$, which yields the rate
\[
\widehat{\xi}_n-\xi_0
=
O_p(h_n^{1/2}).
\]

This generalized framework potentially covers non-DRI drift structures, including polynomially growing drifts, logarithmic trends, almost periodic components, and more general nonstationary deterministic signals.
A systematic investigation of such growth-adapted asymptotic theory will be studied elsewhere.
\end{remark}

\section{Proofs of main theorems}\label{sec:proof}

\subsection{Proof of Theorem \ref{thm:consist}}
First, we shall show that the contrast function \eqref{lg-contrast} converges to a deterministic limit uniformly in $\th\in \ol{\Theta}$. 

By Lemma \ref{lem:f-DX}, we obtain
\begin{align*}
\frac{1}{n} \sum_{i=1}^n  \frac{(\Delta_i^n X - h_n \mu_\xi(t_{i-1}))^2}{2 h_n^2}  \toP - \p_t^2 K(0),
\end{align*}
uniformly in $\xi$. On the other hand, the condition B\ref{as:K2} implies
\begin{align*}
K_\sigma(0) - K_\sigma(h_n) = -\frac{1}{2} \p_t^2 K_\sigma(0) h_n^2 + o(h_n^3),
\end{align*}
uniformly in $\s$. Therefore,
\begin{align*}
\frac{1}{2 [K_\sigma(0) - K_\sigma(h_n)]} = \frac{1}{h_n^2 \p_t^2 K_\sigma(0)} + o(h_n^{-1}), 
\end{align*}
and
\begin{align*}
\log(h_n^{-2}[K_\sigma(0) - K_\sigma(h_n)]) = \log\l( -\frac{1}{2} \p_t^2 K_\sigma(0) \r) + o(h_n),
\end{align*}
uniformly in $\s$. Combining all, we obtain
\begin{align*}
\ell_n(\xi,\s) = \frac{1}{n} \sum_{i=1}^n \l\{ \frac{(\Delta_i^n X - h_n \mu_\xi(t_{i-1}))^2}{-2 h_n^2 \p_t^2 K_\sigma(0)} + \log[-\p_t^2 K_\sigma(0)] \r\} + o_p(1),
\end{align*}
uniformly in $\th$. Since the leading term converges in probability to
\begin{align*}
\ell(\s) := \frac{\p_t^2 K_{\s_0}(0)}{\p_t^2 K_\sigma(0)} + \log[-\p_t^2 K_\sigma(0)]. 
\end{align*}
Hence it follows that 
\[
\sup_{\th \in \ol{\Theta}}\l|\ell_n(\xi,\s) - \ell(\s)\r| \toP 0,\quad n\to \infty. 
\]
Second, note that $\ell(\s)$ is minimized if and only if $\partial_t^2 K_{\sigma_0}(0) = \partial_t^2 K_\sigma(0)$, which implies that $\s=\s_0$ by B\ref{as:id}, that is, it follows that 
\[
\inf_{|\s - \s_0|>\e} |\ell(\s)| > \ell(\s_0). 
\]
Hence, by Theorem 5.7 by van der Vaart \cite{v98}, the following consistency holds true: 
\[
\wh{\s}_n \toP \s_0,\quad n\to \infty. 
\]
Next, note that 
\[
L_n(\xi) = n h_n \l\{ \ell_n(\xi, \wh{\s}_n) - \ell_n(\xi_0, \wh{\s}_n) \r\}. 
\]
Then we write:
\begin{align*}
L_n(\xi) &= h_n\sum_{i=1}^n \l\{ \frac{ h_n^2 \l\{ \mu_{\xi_0}^2(t_{i-1}) - \mu_\xi^2(t_{i-1}) \r\} }{2 [K_{\wh{\s}_n}(0) - K_{\wh{\s}_n}(h_n)]} + \frac{ 2 h_n (\mu_\xi(t_{i-1}) - \mu_{\xi_0}(t_{i-1})) \D_i^n X }{2 [K_{\wh{\s}_n}(0) - K_{\wh{\s}_n}(h_n)]} \r\}.
\end{align*}
By B\ref{as:K2} and $\wh{\s}_n \toP \s_0$,
\begin{align*}
[K_{\wh{\s}_n}(0) - K_{\wh{\s}_n}(h_n)] = -\frac{1}{2} \p_t^2 K(0) h_n^2 + o_p(h_n^3),
\end{align*}
so the reciprocal is
\begin{align*}
\frac{1}{[K_{\wh{\s}_n}(0) - K_{\wh{\s}_n}(h_n)]} = -\frac{2}{\p_t^2 K(0) h_n^2} + o_p(h_n^{-1}).
\end{align*}
Substituting, we obtain
\begin{align*}
L_n(\xi) &= -\frac{1}{\p_t^2 K(0)} \sum_{i=1}^n \l\{ \mu_{\xi_0}^2(t_{i-1}) - \mu_\xi^2(t_{i-1}) \r\}h_n  \\
&\quad - \frac{2}{\p_t^2 K(0)} \sum_{i=1}^n[\mu_\xi(t_{i-1}) - \mu_{\xi_0}(t_{i-1})]\D_i^n X + o_p(h_n).
\end{align*}
Now decompose $\D_i^n X = \D_i^n Z + h_n \mu_{\xi_0}(t_{i-1}) + r_i^n$, where $r_i^n := \int_{t_{i-1}}^{t_i} (\mu_{\xi_0}(s) - \mu_{\xi_0}(t_{i-1})) \df s = o(h_n)$ by continuity.
Then
\begin{align*}
[\mu_\xi(t_{i-1}) - \mu_{\xi_0}(t_{i-1})] \D_i^n X = [\mu_\xi(t_{i-1}) - \mu_{\xi_0}(t_{i-1})] \D_i^n Z + h_n [\mu_\xi(t_{i-1}) - \mu_{\xi_0}(t_{i-1})] \mu_{\xi_0} + o(h_n).
\end{align*}
%Using Lemma \ref{lem:GJ} with the facts that $\E[\D_i^n Z|\F_{t_{i-1}}]=0$ and $\sum_{i=1}^n \E[(\D_i^n Z)^2|\F_{t_{i-1}}]=O_p(nh_n^2)\to 0$, 
It follows from Lemma \ref{lem:m-DZ} with $a_n=h_n^{-1}$ that 
\begin{align*}
\sum_{i=1}^n [\mu_\xi(t_{i-1}) - \mu_{\xi_0}(t_{i-1})]\D_i^n Z \toP 0. 
\end{align*}
Moreover,  by the direct Riemann integrability B\ref{as:DRI-mu} (so $\mu_\xi$ is bounded; Remark \ref{rem:dri-cont}), we have that 
\begin{align*}
\sum_{i=1}^n [\mu_\xi(t_{i-1}) - \mu_{\xi_0}(t_{i-1})] \mu_{\xi_0}(t_{i-1}) h_n \to \int_0^\infty (\mu_\xi(s) - \mu_{\xi_0}(s)) \mu_{\xi_0}(s)\,\df s. 
\end{align*} 
Therefore,
\begin{align*}
L_n(\xi) \to \frac{1}{\p_t^2 K(0)} \int_0^\infty (\mu_\xi(t) - \mu_{\xi_0}(t))^2 \df t =: L(\xi).
\end{align*}
By B\ref{as:id2}, $L(\xi) = 0$ if and only if $\xi = \xi_0$. Hence $L(\xi) > 0$ for all $\xi \ne \xi_0$.
Finally, by Theorem 5.7 by van der Vaart \cite{v98} again, it follows that
\[
\wh{\xi}_n := \arg\min_{\xi \in \ol{\Xi}} \ell_n(\xi, \wh{\s}_n) \toP \xi_0.
\]
This completes the proof. \qed



\subsection{Proof of Theorem \ref{thm:asy-norm}}

Since $\wh{\th}_n\toP \th_0 \in \mathring{\Theta}$, we can assume, in the standard argument for asymptotic normality, 
that $\wh{\th}_n \in \mathring{\Theta}$ for $n$ large enough. without loss of generality. 

Let $D_n$ be the block-diagonal scaling matrix:
\[
D_n :=
\begin{pmatrix}
h_n^{-1/2} I_p & 0 \\
0 & \sqrt{n} I_q
\end{pmatrix}.
\]
Applying Taylor's formula around $\theta_0$, we write
\[
0 = \p_\theta \ell_n(\widehat{\theta}_n)
= \p_\theta \ell_n(\theta_0) + \int_0^1 \p_\theta^2 \ell_n(\theta_n^*(u))\,\df u \cdot (\widehat{\theta}_n - \theta_0),
\]
where $\theta_n^*(u) := u \widehat{\theta}_n + (1 - u) \theta_0$ for some $u \in (0, 1)$.
Multiplying both sides by $D_n$, we obtain
\[
0 = D_n \p_\theta \ell_n(\theta_0) + D_n\int_0^1 \p_\theta^2 \ell_n(\theta_n^*(u))\,\df u \cdot (\widehat{\theta}_n - \theta_0). 
\]
Rewriting, we get
\[
D_n (\widehat{\theta}_n - \theta_0)
= - \l\{C_nD_n \p_\theta^2 \ell_n(\theta_n^*) D_n^{-1} \r\}^{-1} C_nD_n \p_\theta \ell_n(\theta_0), 
\]
where 
\[
C_n :=
\begin{pmatrix}
nh_n I_p & 0 \\
0 &  I_q
\end{pmatrix}.
\]
From Lemma~\ref{lem:score-vector}, we have
\[
C_nD_n \p_\theta \ell_n(\theta_0) \toD \mathcal{N}(0, J(\theta_0)), \quad J(\theta_0) := 
\begin{pmatrix}
\dis  \frac{2\int_0^\infty \l\{ \partial_\xi \mu_{\xi_0}(t) \r\}^{\otimes 2} dt}{\partial_t^2 K(0)} &0 \\
0 & V_1(\s_0)
\end{pmatrix}. 
\]
It remains to show the convergence of $D_n \p_\theta^2 \ell_n(\theta_n^*) D_n^{-1}$. Note that 
\[
C_nD_n \p_\theta^2 \ell_n(\theta_n^*) D_n^{-1}
=
\begin{pmatrix}
nh_n \p_\xi^2 \ell_n(\theta_n^*) & \sqrt{\frac{1}{nh_n}} \p_\xi \p_\sigma^\top \ell_n(\theta_n^*) \\
\sqrt{nh_n} \p_\sigma \p_\xi^\top \ell_n(\theta_n^*) & \p_\sigma^2 \ell_n(\theta_n^*)
\end{pmatrix}.
\]
and we have that 
\begin{align*}
\p_\xi^2 \ell_n(\theta) &= \frac{h_n^2 }{n} \sum_{i=1}^n \l[
\frac{ \l\{ \p_\xi \mu_\xi(t_{i-1}) \r\}^{\otimes 2}}{v_n(\sigma)}
- \frac{ \l\{ \D_i^n X - \mu_\xi(t_{i-1}) h_n \r\} \cdot \p_\xi^2 \mu_\xi(t_{i-1})}{h_n v_n(\sigma)}
\r]; \\
\sqrt{\frac{1}{nh_n}} \p_\xi \p_\sigma \ell_n(\theta)&= \frac{1}{n\sqrt{h_n}} \sum_{i=1}^n
\l[
\frac{ \l\{ \D_i^n X - \mu_\xi(t_{i-1}) h_n \r\} \cdot \p_\xi \mu_\xi(t_{i-1}) h_n \cdot \p_\sigma v_n(\sigma) }
{ \l\{ v_n(\sigma)\r\}^2 }
\r]; \\
\p_\sigma^2 \ell_n(\theta) &= \frac{1}{n} \sum_{i=1}^n \Bigg[
\frac{ \l\{ \D_i^n X - \mu_\xi(t_{i-1}) h_n \r\}^2 \cdot \l[ \p_\sigma^2 \{ v_n(\sigma)\} \cdot \{ v_n(\sigma)\} - 2 \{ \p_\sigma v_n(\sigma) \}^2 \r] }
{ 2 \l\{ v_n(\sigma)\r\}^3 } \\
&+ \frac{ \p_\sigma^2 \{ v_n(\sigma)\} }{ v_n(\sigma)} 
- \frac{ \l\{ \p_\sigma \{ v_n(\sigma)\} \r\}^{\otimes 2} }{ \l\{ v_n(\sigma)\r\}^2 }
\Bigg], 
\end{align*}
where $v_n(\sigma) = K_\sigma(0) - K_\sigma(h_n)$. 

As for $\p_\xi^2 \ell_n(\theta)$, since $\frac{h_n^2}{v_n(\sigma)} = \frac{2}{-\p_t^2 K_\sigma(0)} + o(h_n)$ uniformly in $\s$ by B\ref{as:K2}, 
it follows that 
\[
nh_n\p_\xi^2 \ell_n(\theta) 
= \l( \frac{2}{-\p_t^2 K_\sigma(0)} + o(h_n) \r) \cdot h_n\sum_{i=1}^n \l[\l\{ \p_\xi \mu_\xi(t_{i-1}) \r\}^{\otimes 2} - \widetilde{Y}_i^n \cdot \p_\xi^2 \mu_\xi(t_{i-1})\r],
\]
where $\widetilde{Y}_i^n := \frac{ \Delta_i^n X - h_n \mu_\xi(t_{i-1}) }{h_n}$.
Now, by the assumption B\ref{as:DRI-mu2}, the first term of the summation satisfies
\[
h_n\sum_{i=1}^n \l\{ \p_\xi \mu_\xi(t_{i-1}) \r\}^{\otimes 2}  =  \sum_{i=1}^n \l\{ \p_\xi \mu_\xi(t_{i-1}) \r\}^{\otimes 2}h_n   
\to \int_0^\infty \{ \p_\xi \mu_\xi(s) \}^{\otimes 2}\df s, 
\]
as $n\to \infty$ since $(\p_\xi \mu_\xi)^2$ is DRI. Moreover, it follows for the second term that 
\begin{align*}
h_n\sum_{i=1}^n \widetilde{Y}_i^n \cdot \p_\xi^2 \mu_\xi(t_{i-1}) 
&=\sum_{i=1}^n\l[\D_i^n Z + \int_{t_{i-1}}^{t_i}\mu_{\xi_0}(s)\,\df s - h_n \mu_\xi(t_{i-1})  \r]  \cdot \p_\xi^2 \mu_\xi(t_{i-1}) \\
&= \sum_{i=1}^n \p_\xi^2 \mu_\xi(t_{i-1})  \cdot \D_i^nZ 
+ h_n \sum_{i=1}^n \frac{1}{h_n}\int_{t_{i-1}}^{t_i}\mu_{\xi_0}(s)\,\df s  \cdot \p_\xi^2 \mu_\xi(t_{i-1}) \\
&\quad -  \sum_{i=1}^n \mu_\xi(t_{i-1})  \cdot \p_\xi^2 \mu_\xi(t_{i-1}) h_n \\
&\to  \int_0^\infty \mu_{\xi_0}(s) \p_\xi^2 \mu_\xi(s) \df s - \int_0^\infty \mu_\xi(s) \p_\xi^2 \mu_\xi(s) \df s, 
\end{align*}
uniformly in $\xi \in \Xi$ by Lemma \ref{lem:m-DZ} with $a_n=h_n^{-1}$ and the mean value theorem under B\ref{as:DRI-mu2} and B\ref{as:DRI-mu4}. Therefore,
\[
nh_n\p_\xi^2 \ell_n(\theta) \to \frac{2}{\p_t^2 K_\sigma(0)} \int_0^\infty \l[- \mu_{\xi_0}(s) \p_\xi^2 \mu_\xi(s)  + \mu_\xi(s) \p_\xi^2 \mu_\xi(s) -  \{ \p_\xi \mu_\xi(s) \}^{\otimes 2}\r] \df s, 
\]
uniformly in $\theta \in \overline{\Theta}$. Hence we have that 
\[
nh_n \p_\xi^2 \ell_n(\wh{\theta}_n) \to - \frac{2}{\p_t^2 K_{\sigma_0}(0)}  \int_0^\infty \{ \p_\xi \mu_\xi(s) \}^{\otimes 2}\,\df s.
\]
As for $\p_\sigma^2 \ell_n(\theta)$, it follow by the same argument as above that 
\begin{align*}
\partial_\sigma^2 \ell_n(\theta)
&= \frac{1}{n} \sum_{i=1}^n \l[
\frac{ \widetilde{Y}_i^n{}^2 \cdot \l( \partial_\sigma^2 \partial_t^2 K_\sigma(0) \cdot \partial_t^2 K_\sigma(0) 
- \l( \partial_\sigma \partial_t^2 K_\sigma(0) \r)^{\otimes 2} \r) }
{ \l( \partial_t^2 K_\sigma(0) \r)^3 } + o(1)
\r] \\
&= \l( \frac{ \partial_\sigma^2 \partial_t^2 K_\sigma(0) \cdot \partial_t^2 K_\sigma(0)
- \l( \partial_\sigma \partial_t^2 K_\sigma(0) \r)^{\otimes 2} }{ \l( \partial_t^2 K_\sigma(0) \r)^3 } \r)
\cdot \frac{1}{n} \sum_{i=1}^n \widetilde{Y}_i^n{}^2 + o_p(1).
\end{align*}

By Lemma~\ref{lem:f-DX} with $f(z, \theta) = z^2$, we have:
\[
\frac{1}{n} \sum_{i=1}^n \widetilde{Y}_i^n{}^2 \to \E[Z^2] = -\partial_t^2 K(0),
\]
and thus
\[
\partial_\sigma^2 \ell_n(\theta)
\to 
\l( \frac{ \partial_\sigma^2 \partial_t^2 K_\sigma(0) \cdot \partial_t^2 K_\sigma(0)
- \l( \partial_\sigma \partial_t^2 K_\sigma(0) \r)^{\otimes 2} }{ \l( \partial_t^2 K_\sigma(0) \r)^3 } \r)  (-\partial_t^2 K_{\sigma_0}(0)).
\]
uniformly in $\th\in \Theta$. Hence, we have that 
\begin{align*}
\partial_\sigma^2 \ell_n(\wh{\theta}_n) &\toP 
\frac{ \partial_\sigma^2 \partial_t^2 K_{\sigma_0}(0) }{ \partial_t^2 K_{\sigma_0}(0) }
- \l( \frac{ \partial_\sigma \partial_t^2 K_{\sigma_0}(0) }{ \partial_t^2 K_{\sigma_0}(0) } \r)^{\otimes 2} 
= \p_\s^2\l(\log (-\p_t^2 K_{\s_0}(0) )\r) = V_2(\s_0). 
\end{align*}
Similarly, we see that $ \frac{1}{\sqrt{nh_n}} \p_\xi \p_\sigma \ell_n(\theta)=o_p(1)$ and $\sqrt{nh_n} \p_\xi \p_\sigma \ell_n(\theta)=o_p(1)$.  Hence 
\[
D_n \p_\theta^2 \ell_n(\theta_n^*) D_n^{-1}
\toP I(\th_0):=
\begin{pmatrix}
\dis \frac{2}{\p_t^2 K_{\sigma_0}(0)}  \int_0^\infty \{ \p_\xi \mu_\xi(s) \}^{\otimes 2}\,\df s& 0 \\
0 &V_2(\s_0)
\end{pmatrix}.
\]

As a consequence, we have 
\[
D_n (\widehat{\theta}_n - \theta_0)\toD {\cal N}\l(\mb{0}, I^{-1}(\th_0)J(\th_0) I^{-1}(\th_0)\r) = \mathcal{N}(0, \Sigma(\theta_0)). 
 \]
\qed

\subsection{Proof of Theorem \ref{thm:M-consist}}

We may apply Lemma~\ref{lem:fY}, which yields the uniform convergence
\[
\sup_{\sigma \in \overline{\Pi}} \l| \Phi_n(\sigma) - \Phi(\sigma) \r| \toP 0.
\]
In addition, since $\wt{\s}_n$ satisfies $\Phi_n(\wt{\s}_n) = 0$, and since $\Phi$ satisfies the identifiability condition \eqref{Phi-identifiable}, it follows from standard $Z$-estimation theory: 
\[
\wt{\s}_n \toP \sigma_0.
\]
See, e.g., Theorem 5.9 in van der Vaart \cite{v98}. 
\qed


\subsection{Proof of Theorem \ref{thm:M-asy-norm}}
By definition, the estimator $\wt{\s}_n \in \mathbb{R}^q$ satisfies
\[
\Phi_n(\wt{\s}_n) = \mb{0}.
\]
According to (the integral form of) the mean value theorem, we obtain
\[
\Phi_n(\wt{\s}_n) - \Phi_n(\sigma_0) 
= \l( \int_0^1 \partial_\sigma \Phi_n\l( \sigma_0 + u (\wt{\s}_n - \sigma_0) \r) \mathrm{d}u \r) (\wt{\s}_n - \sigma_0).
\]
Hence,
\[
\sqrt{n} (\wt{\s}_n - \sigma_0)
= - \l( \int_0^1 \partial_\sigma \Phi_n\l( \sigma_0 + u (\wt{\s}_n - \sigma_0) \r) \mathrm{d}u \r)^{-1} 
\sqrt{n} \Phi_n(\sigma_0).
\]

Since $\wt{\s}_n \to \sigma_0$ in probability and $\partial_\sigma \Phi_n(\sigma)$ converges uniformly in probability to $\partial_\sigma \Phi(\sigma)$ on a neighborhood of $\sigma_0$, we obtain
\[
\int_0^1 \partial_\sigma \Phi_n\l( \sigma_0 + u (\wt{\s}_n - \sigma_0) \r) \mathrm{d}u 
\toP \partial_\sigma \Phi(\sigma_0) = A,
\]
and the inverse converges in probability to $A^{-1}$.

Next, we apply Theorem 2.1 of Neumann~\cite{n13} to obtain that 
\[
\sqrt{n} \Phi_n(\sigma_0) \toD \mathcal{N}(\mb{0}, \Gamma^2),\quad n\to \infty, 
\]
with  $\dis \Gamma^2 := \lim_{n \to \infty} n \, \mathrm{Var}\l( \Phi_n(\sigma_0) \r).$

Note that
\[
\Phi_n^{(j)}(\sigma_0) = \frac{1}{n} \sum_{i=1}^n \l[ f_j(Y_{i-1}^n) 
- \int_{\mathbb{R}} f_j(z) \, \phi_{K_{\sigma_0}(0)}(z) \, \mathrm{d}z \r]
=: \frac{1}{n} \sum_{i=1}^n\zeta_{i,n}^{(j)}.
\]
Since $Y_i^n = Z_{t_i} + \Delta_i^n$, where 
\[
\Delta_i^n := \int_0^{t_i} \l\{ \mu_{\xi_0}(s) - \mu_{\widehat{\xi}_n}(s) \r\} \mathrm{d}s,
\]
we define the centered version
\[
\widetilde{\zeta}_{i,n}^{(j)} := f_j(Y_{i-1}^n) - \E[f_j(Y_{i-1}^n)],
\]
so that $\E[\widetilde{\zeta}_{i,n}^{(j)} ] = 0$, and write
\[
\Phi_n^{(j)}(\sigma_0)
= \frac{1}{n} \sum_{i=1}^n \widetilde{\zeta}_{i,n}^{(j)} + R_n^{(j)},
\]
where
\[
R_n^{(j)} := \frac{1}{n} \sum_{i=1}^n \l\{ \E[f_j(Y_{i-1}^n)] - \int_{\mathbb{R}} f_j(z) \phi_{K_{\sigma_0}(0)}(z) \mathrm{d}z \r\}.
\]

By a Taylor expansion and the consistency $\widehat{\xi}_n \to \xi_0$, we have $\sup_i |\Delta_i^n| = o_P(1)$ and hence
\[
R_n^{(j)} = o_P(n^{-1/2}).
\]

To apply Theorem 2.1 of Neumann~\cite{n13} to $\{ \widetilde{\zeta}_{i,n}^{(j)} \}_{i=1}^n$, we verify the following:

\begin{itemize}
\item[(i)] \textbf{Square integrability:}
There exists a constant $v_0 > 0$ such that
\[
\sum_{i=1}^n \E\l[\, (\widetilde{\zeta}_{i,n}^{(j)} )^2\, \r] \le v_0
\quad \text{for all } n \in \mathbb{N}.
\]
This follows from the polynomial growth of $f_j$ and bounded moments of $Z_{t_i}$, together with the fact that $Y_i^n = Z_{t_i} + \Delta_i^n$ and $\Delta_i^n = o_P(1)$.

\item[(ii)] \textbf{Lindeberg-type condition:}
For all $\varepsilon > 0$,
\[
\sum_{i=1}^n \E\l[\, (\widetilde{\zeta}_{i,n}^{(j)} )^2 \cdot \mathbf{1}_{\{ |\widetilde{\zeta}_{i,n}^{(j)} | > \varepsilon \sqrt{n} \}} \, \r] \to 0 \quad (n \to \infty).
\]
This is ensured by the same growth and moment conditions as in (ii), combined with the fact that $Z_{t_i}$ is Gaussian.

\item[(iii)] \textbf{Weak dependence (covariance inequality):}
There exists a sequence $\{\theta_r\}_{r\in \mathbb{N}}$ with $\sum_{r=1}^\infty \theta_r < \infty$ such that, for any measurable function $g : \mathbb{R}^u \to \mathbb{R}$ with $|g| \le 1$,
\[
\l| \Cov\l( g(\widetilde{\zeta}_{s_1,n}^{(j)},\dots,\widetilde{\zeta}_{s_u,n}^{(j)}), \widetilde{\zeta}_{t,n}^{(j)} \r) \r|
\le ( \E[\widetilde{\zeta}_{t,n}^{(j)2}] + \E[\widetilde{\zeta}_{s_u,n}^{(j)2}] + n^{-1} ) \cdot \theta_{r},
\]
where $s_1 < \cdots < s_u < s_u + r \le t$.
\item[(iv)] \textbf{Convergence of variance:}
There exists a constant $\Gamma_{j}^2 \in (0,\infty)$ such that
\[
\Var\l( \sum_{i=1}^n \widetilde{\zeta}_{i,n}^{(j)} \r) \to \Gamma_{j}^2.
\]
\end{itemize}

\paragraph{Verification of (i) Square integrability.}
Recall that $\widetilde{\zeta}_{i,n}^{(j)} = f_j(Y_{i-1}^n) - \E[f_j(Y_{i-1}^n)]$, where $Y_i^n = Z_{t_i} + \delta_i^n$ with 
\[
\delta_i^n := \int_0^{t_i} \{ \mu_{\xi_0}(s) - \mu_{\widehat{\xi}_n}(s) \} \, \mathrm{d}s.
\]
By Jensen's inequality,
\[
\E\l[ (\widetilde{\zeta}_{i,n}^{(j)} )^2 \r] \le \E\l[ f_j(Y_{i-1}^n)^2 \r].
\]
Under assumption~(B-1), there exists $C > 0$ such that $|f_j(x)| \lesssim 1 + |x|^C$, hence
\[
f_j(Y_{i-1}^n)^2 \lesssim 1 + |Z_{t_i} + \delta_i^n|^{2C} \lesssim 1 + |Z_{t_i}|^{2C} + |\delta_i^n|^{2C}.
\]
Since $\{Z_{t_i}\}_{i=1}^n$ has uniformly bounded moments and $\sup_i |\delta_i^n| = o_P(1)$, we obtain
\[
\sup_{n} \sum_{i=1}^n \E\l[ (\widetilde{\zeta}_{i,n}^{(j)} )^2 \r] < \infty.
\]
Therefore, condition~(i) is satisfied.

\paragraph{Verification of (ii) Lindeberg-type condition.}
Let $\varepsilon > 0$ be arbitrary. By assumption~(B-1), we have $|f_j(x)| \lesssim 1 + |x|^C$ for some $C > 0$, so
\[
|\widetilde{\zeta}_{i,n}^{(j)} | \lesssim 1 + |Y_{i-1}^n|^C, \quad Y_i^n = Z_{t_i} + \delta_i^n.
\]
Since $\delta_i^n = o_P(1)$ uniformly and $Z_{t_i}$ has uniformly bounded moments of all orders, we obtain for any $r > 2$
\[
\sup_{n,i} \E\l[ |\widetilde{\zeta}_{i,n}^{(j)} |^r \r] < \infty.
\]
Then it follows from Markov's inequality that 
\[
\sum_{i=1}^n \E\l[ (\widetilde{\zeta}_{i,n}^{(j)} )^2 \cdot \mathbf{1}_{\{ |\widetilde{\zeta}_{i,n}^{(j)} | > \varepsilon \sqrt{n} \}} \r]
\lesssim n \cdot \frac{1}{n^{(r-2)/2}} \to 0 \quad (n \to \infty).
\]
Thus, the Lindeberg-type condition~(ii) is satisfied.


\paragraph{Verification of (iii) Weak dependence.}
Fix integers $1 \le s_1 < \dots < s_u < t \le n$ and define $r := t - s_u$. Let $g : \mathbb{R}^u \to \mathbb{R}$ be any measurable function with $|g| \le 1$. Set
\[
\theta_r := \frac{ \displaystyle \int_0^{h_n} \int_0^{h_n} |K(r h_n + u - v)| \,\df u \,\df v }{
\displaystyle 2 \int_0^{h_n} \int_0^{h_n} |K(u - v)| \,\df u \,\df v + h_n }.
\]

We write $Y_i^n = Z_{t_i} + \delta_i^n$ and note that $\delta_i^n = o_P(1)$ uniformly in $i$. Since $f_j$ is of polynomial growth and differentiable, we can linearize $f_j(Y_{i-1}^n)$ around $Z_{t_i}$, yielding
\[
f_j(Y_{i-1}^n) = f_j(Z_{t_{i-1}}) + R_i^n,
\]
where the remainder $R_i^n = f_j(Z_{t_{i-1}} + \delta_i^n) - f_j(Z_{t_{i-1}})$ satisfies $R_i^n = o_P(1)$ under the growth condition.

Let us define the auxiliary array
\[
X_i^n := \frac{1}{\sqrt{n}} f_j(Z_{t_i}), \quad \text{so that } \widetilde{\zeta}_{i,n}^{(j)} = X_i^n - \E[X_i^n] + o_P(n^{-1/2}).
\]
By Lemma~\ref{lem:cov-bound-summable} (adapted to this setting), for such triangular arrays we have
\[
\l| \Cov\l( g(X_{s_1}^n,\dots,X_{s_u}^n), X_t^n \r) \r| 
\le \l( \E[|X_t^n|^2] + \E[|X_{s_u}^n|^2] + \frac{1}{n} \r) \cdot \theta_r.
\]

Since the centering operation and the $o_P(n^{-1/2})$ term do not affect the covariance up to $o(n^{-1})$, it follows that
\[
\l| \Cov\l( g(\widetilde{\zeta}_{s_1,n}^{(j)}, \dots, \widetilde{\zeta}_{s_u,n}^{(j)}), \widetilde{\zeta}_{t,n}^{(j)} \r) \r| 
\le \l( \E[|\widetilde{\zeta}_{t,n}^{(j)}|^2] + \E[|\widetilde{\zeta}_{s_u,n}^{(j)}|^2] + \frac{1}{n} \r) \cdot \theta_r + o(n^{-1}).
\]

Finally, if $K \in L^1(\mathbb{R})$, then $\sum_{r=1}^\infty \theta_r < \infty$ by Lemma~\ref{lem:cov-bound-summable}, and thus the weak dependence condition (iii) is satisfied.

\paragraph{Verification of (iv) Convergence of variance.}
Recall the definition:
\[
\widetilde{\zeta}_{i,n}^{(j)} := \frac{1}{\sqrt{n}} \l\{ f_j(Y_{i-1}^n) - \E[f_j(Y_{i-1}^n)] \r\}, \quad Y_i^n := Z_{t_i} + \delta_i^n,
\]
with $\delta_i^n := -\int_0^{t_i} \mu_{\widehat{\xi}_n}(s) \,\df s + \int_0^{t_i} \mu_{\xi_0}(s) \,\df s = o_P(1)$ uniformly in $i$.

Let us define $X_i^n := \frac{1}{\sqrt{n}} \l\{ f_j(Z_{t_{i-1}}) - \E[f_j(Z_{t_{i-1}})] \r\}$. Then we have
\[
\widetilde{\zeta}_{i,n}^{(j)} = X_i^n + R_i^n, \quad \text{with } R_i^n := \frac{1}{\sqrt{n}} \l\{ f_j(Y_{i-1}^n) - f_j(Z_{t_{i-1}}) - \E[f_j(Y_{i-1}^n) - f_j(Z_{t_{i-1}})] \r\}.
\]

By the smoothness and polynomial growth of $f_j$, and the fact that $\delta_i^n = o_P(1)$ uniformly, we obtain
\[
\sup_{1 \le i \le n} |R_i^n| = o_P(n^{-1/2}),
\quad \text{hence } \sum_{i=1}^n R_i^n = o_P(1).
\]

It follows that
\[
\sum_{i=1}^n \widetilde{\zeta}_{i,n}^{(j)} = \sum_{i=1}^n X_i^n + o_P(1),
\quad \text{so that } \Var\l( \sum_{i=1}^n \widetilde{\zeta}_{i,n}^{(j)} \r)
= \Var\l( \sum_{i=1}^n X_i^n \r) + o(1).
\]

Now consider
\[
\Var\l( \sum_{i=1}^n X_i^n \r)
= \sum_{i=1}^n \Var(X_i^n) + 2 \sum_{1 \le i < j \le n} \Cov(X_i^n, X_j^n).
\]
Since $X_i^n = \frac{1}{\sqrt{n}} \{ f_j(Z_{t_i}) - \E[f_j(Z_{t_i})] \}$ and $\{ Z_{t_i} \}$ is a stationary Gaussian process, the sequence $f_j(Z_{t_i})$ is stationary and $\alpha$-mixing under appropriate conditions on $K$.
Therefore, by standard results for weakly dependent stationary sequences (see e.g., Ibragimov and Rozanov~(1978), Doukhan~(1994)), we have the convergence
\[
\Var\l( \sum_{i=1}^n X_i^n \r)
= \sum_{h = -(n-1)}^{n-1} \l(1 - \frac{|h|}{n}\r) \Cov\l( f_j(Z_0), f_j(Z_{|h|}) \r) \to \Gamma_{j}^2,
\]
where
\[
\Gamma_{j}^2 := \sum_{h \in \mathbb{Z}} \Cov\l( f_j(Z_0), f_j(Z_h) \r) \in (0,\infty),
\]
provided that $K \in L^1(\mathbb{R})$ and $f_j$ is of polynomial growth.
Combining the above with $\sum_{i=1}^n R_i^n = o_P(1)$, we conclude
\[
\Var\l( \sum_{i=1}^n \widetilde{\zeta}_{i,n}^{(j)} \r) \to \Gamma_{j}^2,
\]
as required. Therefore, by Theorem 2.1 of Neumann~\cite{n13},
\[
\sqrt{n} \Phi_n^{(j)}(\sigma_0) \xrightarrow{d} \mathcal{N}(0, \Gamma_{j}^2),
\]
for some variance $\Gamma_{j}^2$. Applying this component-wise for $j=1,\dots,q$, we obtain the vectorial convergence
\[
\sqrt{n} \Phi_n(\sigma_0) \xrightarrow{d} \mathcal{N}(\mb{0}, \Gamma^2).
\]
Combining with the convergence of $A_n^{-1} \toP A^{-1}$, we conclude that
\[
\sqrt{n}(\wt{\s}_n - \sigma_0) \xrightarrow{d} \mathcal{N}(0, A^{-1} \Gamma^2 A^{-1}),
\]
as desired.


\qed



\ \vspace{1cm}\ \\
\noindent {\bf\large Acknowledgement.} 
%The authors sincerely thank the anonymous reviewers for their insightful comments, which have enhanced the quality of this paper. 
This work is partially supported by JSPS KAKENHI Grant-in-Aid for Scientific Research (C) \#24K06875; Japan Science and Technology Agency CREST \#JPMJCR2115. 

\
\appendix
\section{Limit theorems for stationary Gaussian processes}\label{app:A}

\begin{lem}\label{lem:m-DZ}
Let $Z=(Z_t)_{t\ge0}\sim GP(0,K)$ be centered and stationary.
Assume A\ref{as:local-kernel} and A\ref{as:increment-covariance}.
Let $m:[0,\infty)\to\mathbb{R}$ be measurable and suppose that
\[
\frac1{a_n}\sum_{i=1}^n m(t_{i-1})^2=O(1).
\]
Then
\[
\sum_{i=1}^n m(t_{i-1})\Delta_i^nZ=O_P\left((a_nh_n^2)^{1/2}\right).
\]
\end{lem}

\begin{proof}
Set $S_n(m):=\sum_{i=1}^n m(t_{i-1})\Delta_i^nZ$.
Since $Z$ is centered, we have $\E[S_n(m)]=0$.
Moreover,
\[
\Var(S_n(m))=\sum_{i,j=1}^n m(t_{i-1})m(t_{j-1})\Cov(\Delta_i^nZ,\Delta_j^nZ).
\]
Using the inequaity $2|xy|\le x^2+y^2$, we obtain
\begin{align*}
|\Var(S_n(m))|&\le \frac12\sum_{i,j=1}^n \{m(t_{i-1})^2+m(t_{j-1})^2\}\left|\Cov(\Delta_i^nZ,\Delta_j^nZ)\right| \\
&\le \sum_{i=1}^n m(t_{i-1})^2\sum_{j=1}^n \left|\Cov(\Delta_i^nZ,\Delta_j^nZ)\right|. 
\end{align*}
By A\ref{as:increment-covariance},
\[
|\Var(S_n(m))|\le C h_n^\beta\sum_{i=1}^n m(t_{i-1})^2.
\]
Since $a_n^{-1}\sum_{i=1}^n m(t_{i-1})^2=O(1)$, we get $\Var(S_n(m))=O(nh_n^2$.
Therefore, $S_n(m)=O_P((a_nh_n^2)^{1/2})$.
\end{proof}



\begin{lem}\label{lem:diff-L2}
Let $Z = (Z_t)_{t \ge 0} \sim GP(0,K)$ with the conditions A\ref{as:2} and A\ref{as:3}. 
Then, there eixists the `mean-square derivative' $\dot{Z}_t$ in the sense that 
\begin{align}
\frac{Z_{t+h} - Z_t}{h} \stackrel{L^2}{\longrightarrow} \dot{Z}_t,\quad h\to 0, \label{diff-L2}
\end{align}
and $\dot{Z} \sim GP\l(0, -\p^2_t K(0)\r)$. In addition, suppose A\ref{as:4}. Then, $\dot{Z}$ is ergodic in the sense of Corollary  \ref{cor:ergodic}.
\end{lem}

\begin{proof}
See Ibragimov and Rozanov \cite{ir78}, Section I.7.1.
\end{proof}

\begin{lem}\label{lem:approx-integral}
Let $Z = (Z_t)_{t \ge 0} \sim GP(0,K)$ with the conditions A\ref{as:2}--A\ref{as:4}. 
Then, for all functions $f:\mathbb{R} \to \mathbb{R}$ which is continuous and is of polynomial growth:
\[
|f(x)| \le C(1 + |x|^p), \quad \text{for all } x \in \mathbb{R},
\]
for some $C>0$ and $p \ge 1$, it holds that 
\[
\frac{1}{nh_n} \int_0^{nh_n} \l\{ f\l( \frac{Z_{t+h_n} - Z_t}{h_n} \r) - f(\dot{Z}_t) \r\} dt \stackrel{L^1}{\longrightarrow} 0, 
\]
under $h_n \to 0$ and $nh_n \to \infty$ as $n\to \infty$. 
\end{lem}

\begin{proof}
Since $Z$ is mean-square differentiable \eqref{diff-L2}, also in probability,  it holds that 
\[
f\l( \frac{Z_{t+h_n} - Z_t}{h_n} \r) \toP f(\dot{Z}_t),\quad n\to \infty, 
\]
by the continuous mapping theorem. 

To conclude $L^1$ convergence of the integrals, we verify uniform integrability. By the growth condition on $f$, we have
\[
\l| f\l( \frac{Z_{t+h_n} - Z_t}{h_n} \r) \r| \le C(1 + |Y_n(t)|^p),
\]
where $Y_n(t) := \frac{Z_{t+h_n} - Z_t}{h_n}$. We estimate
\[
\E\l[ |f(Y_n(t))| \cdot \mathbf{1}_{\{ |f(Y_n(t))| > K \}} \r]
\le C \E\l[ (1 + |Y_n(t)|^p) \cdot \mathbf{1}_{\{ |Y_n(t)| > \lambda \}} \r],
\]
for $\lambda := \l( \frac{K}{C} - 1 \r)^{1/p}$. By H\"older's inequality,
\[
\E[|Y_n(t)|^p \cdot \mathbf{1}_{\{ |Y_n(t)| > \lambda \}}]
\le \l( \E[|Y_n(t)|^q] \r)^{p/q} \cdot \l( \mathbb{P}(|Y_n(t)| > \lambda) \r)^{1 - p/q}, 
\]
for $p>0$ and $q\in (1,2)$ with $1/p + 1/q=1$. 
Since $Y_n(t) \to \dot{Z}_t$ in $L^2$, the sequence $\{Y_n(t)\}$ is bounded in $L^q$ uniformly in $n$, and the tail probability decays rapidly since it has a Gaussian tail. 
Thus the upper bound in the last right-hand side tends to $0$ as $\lambda \to \infty$, uniformly in $n$. Therefore, the family $\{ f(Y_n(t)) \}_n$ is uniformly integrable.
Hence it follows by Vitali's convergence theorem that 
\[
\E\l| f\l( \frac{Z_{t+h_n} - Z_t}{h_n} \r) - f(\dot{Z}_t) \r|\to  0,\quad n\to \infty. 
\]
Finally, using Fubini's theorem and dominated convergence,
\begin{align*}
&\E \l|\frac{1}{nh_n} \int_0^{nh_n}\l[ f\l( \frac{Z_{t+h_n} - Z_t}{h_n} \r) - f(\dot{Z}_t) \r] \df t \r| \\
&= \frac{1}{nh_n} \int_0^{nh_n}\E \l| f\l( \frac{Z_{t+h_n} - Z_t}{h_n} \r) - f(\dot{Z}_t) \r|\,\df t \to 0,\quad n\to \infty. 
\end{align*}
This completes the proof. 
\end{proof}

\begin{lem}\label{lem:avg-conv}
Let $Z = (Z_t)_{t \ge 0} \sim GP(0,K)$ with A\ref{as:2}--A\ref{as:4}, and let $f:\mathbb{R}\times \ol{\Theta} \to \mathbb{R}$ be continuous and of polynomial growth:
\begin{align*}
|\p_\th^k f(x,\th)| \le C(1 + |x|^p), \quad \text{for all } x \in \mathbb{R},\ k=0,1, %label{lem:avg-conv}
\end{align*}
for some $C>0$ and $p \ge 1$. Then, it holds that 
\[
\frac{1}{n} \sum_{i=1}^n f\l( \frac{\Delta_i^n Z}{h_n} ,\th \r) \toP \int_\R f(z,\th)\phi_{-\p^2_t K(0)}(z)\,\df z, 
\]
uniformly in $\th\in \ol{\Theta}$, under $h_n \to 0$, $n h_n \to \infty$ as $n \to \infty$. 
\end{lem}

\begin{proof}
First, we shall show the convergence for each fixed $\th\in \ol{\Theta}$. 

Accroding to the stationarity of $Z$, we see that 
\begin{align*}
Y_i^n := \frac{\Delta_i^n Z}{h_n} = h_n^{-1} (1 , -1)\begin{pmatrix}Z_{t_i^n} \\ Z_{t_{i-1}^n} \end{pmatrix} \sim N\l(0, 2[K(0) - K(h_n)]\r). 
\end{align*}
Hence 
\begin{align*}
Y_i^n \toD \dot{Z}_0 \sim {\cal N}(0, -\partial^2 K(0)),\quad n\to \infty. 
\end{align*}
Therefore, we shall show that, for each $\th\in \ol{\Theta}$, 
\[
G_n(\th):=\frac{1}{n}\sum_{i=1}^n f(Y_{i-1}^n,\th)  - \E[f(\dot{Z}_0,\th)] \toP 0,\quad n\to \infty.
\]
Since $Y_{i-1}^n  \sim N\l(0, 2[K(0) - K(h_n)]\r)$ and $\dot{Z}_0 \sim N\l(0,-\p^2_t K(0)\r)$, it follows by the dominated convergence theorem that 
\[
\lim_{n \to \infty} \E[f(Y_{i-1}^n,\th)] = \E[f(\dot{Z}_0,\th)] = \int_\R f(z,\th)\phi_{-\p^2_t K(0)}(z)\,\df z.
\]
Define $g_\th(t,h) := f\l( \frac{Z_{t+h} - Z_t}{h} \r)$, so that $f(Y_{i-1}^n,\th) = g_\th(t_{i-1},h_n)$. 
Then, using the notation in Lemma \ref{lem:approx-integral}, we have
\begin{align*}
\l|G_n(\th)\r| &\le \l|\frac{1}{nh_n}\sum_{i=1}^n g_\th(t_{i-1},h_n) h_n - \frac{1}{nh_n}\int_0^{nh_n} g_\th(t,h_n)\,\df t \r|\\
&\quad + \l|\frac{1}{nh_n}\int_0^{nh_n} \l[ g_\th(t,h_n) - f(\dot{Z}_t,\th) \r]\,\df t \r| + \l|\frac{1}{nh_n}\int_0^{nh_n}f(\dot{Z}_t,\th)\,\df t -  \E[f(\dot{Z}_0,\th)]\r|\\
& =: G_n^1 + G_n^2+ G_n^3. 
\end{align*}
Note that $G_n^2\stackrel{L^1}{\longrightarrow} 0$ by Lemma \ref{lem:approx-integral}, and that $G_n^3 \to 0\ a.s.\ (n\to \infty)$ by the ergodicity of $\dot{Z}$.  
To complete the proof, we show $G_n^1 \to 0\ a.s.$ as $n \to \infty$. 

Note that, for fixed $h > 0$, the function $t \mapsto g_\th(t,h)$ is continuous in $t$, since $Z$ has continuous sample paths almost surely and $f$ is continuous.
On each subinterval $[t_{i-1}, t_i]$, by the mean value theorem for integrals, there exists $\tau_i \in [t_{i-1}, t_i]$ such that
\[
\int_{t_{i-1}}^{t_i} g_\th(t, h_n) \,\df t = g_\th(\tau_i, h_n) h_n\quad a.s.
\]
Therefore, we can write
\[
G_n^1 = \l| \frac{1}{nh_n} \sum_{i=1}^n \l[ g_\th(t_{i-1}, h_n) - g_\th(\tau_i, h_n) \r] h_n \r| 
= \frac{1}{n} \sum_{i=1}^n \l| g_\th(t_{i-1}, h_n) - g_\th(\tau_i, h_n) \r|.
\]
Since $g(t,h)$ is uniformly continuous in $t$ on compacts, and $|\tau_i - t_{i-1}| \le h_n \to 0$ we see that 
\[
G_n^1 \le \sup_{|s - t| \le h_n} |g_\th(s, h_n) - g_\th(t, h_n)| \to 0\quad a.s., 
\]
as $h_n\to 0$. 

For the uniformity of convergence, we shall show the tightness of the random functions $G_n(\th) (n=1,2,\dots)$, which is confirmed via the following tightness criterion: 
\[
\sup_n\E\l[ \sup_{\th\in \ol{\Theta}} \l|\p_\th G_n(\th)\r|\r] < \infty. 
\]
Actually, it is easy to see by the condition \eqref{lem:avg-conv} that 
\begin{align*}
\sup_{\th\in \ol{\Theta}} \l|\p_\th G_n(\th)\r| \lesssim \frac{1}{n}\sum_{i=1}^n (1 + |Y_{i-1}^n|^C) + \int_\R (1 + |z|^C)\phi_{-\p^2_t K(0)}(z)\,\df z, 
\end{align*}
which is integrable uniformly in $n\in \N$. Hence, the proof is completed. 
\end{proof}


\begin{cor}\label{cor:moment}
Under the same assumptions as in Lemma \ref{lem:avg-conv}, it follows for any integer $\k\ge 1$ that
\[
\frac{1}{nh_n^ {2\k}} \sum_{i=1}^n (\D_i^n Z)^{2\k} \toP \frac{(2\k)!}{2^\k\cdot \k!} \l(-\p^2_t K(0)\r)^\k, 
\]
as $n\to \infty$. 
\end{cor}

\begin{proof}
Take $f(x) = x^\k$ in Lemma \ref{lem:avg-conv}, and note that $\E[G^\k] = \frac{(2\k)!}{2^\k\cdot \k!} \s^{2\k}$ for $G\sim {\cal N}(0,\s^2)$. 

\end{proof}


\begin{lem}\label{lem:ergod}
Let $Z = (Z_t)_{t \in \mathbb{R}}$ be a centered stationary Gaussian process with the conditions A\ref{as:2} and  A\ref{as:3}, and 
let $f(x,\th): \R\times \ol{\Theta} \to \R$ be a measureble function such that 
\begin{align}
\sup_{\th \in \ol{\Theta}}|\p_x^k \p_\th^l f(x,\th)| \lesssim 1 + |x|^C, \label{cond:ergod}
\end{align}
for integers $k$ and $l$ with $0\le k+l\le 1$. 
Then, under the conditions that $h_n \to 0$ and $nh_n\to \infty$ as $n\to \infty$, the following hold true:
\[
\sup_{\th \in \ol{\Theta}}\l|\frac{1}{n} \sum_{i=1}^n f(Z_{t_{i-1}^n},\th) - \int_\R f(z,\th)\phi_{K(0)}(z)\,\df z\r|\toP 0,
\]
as $n\to \infty$. 
\end{lem}

\begin{proof}
Note that
\begin{align*}
&\E\l| \frac{1}{n} \sum_{i=1}^n f(Z_{t_{i-1}^n},\th)  - \frac{1}{nh_n} \int_0^{nh_n} f(Z_u,\th)\,\df u \r| \\
&\le \frac{1}{nh_n} \sum_{i=1}^n \E\int_{t_{i-1}^n}^{t_i^n} |f(Z_{t_{i-1}^n},\th) - f(Z_u,\th)|\,\df u \\
&\le  \frac{1}{nh_n} \sum_{i=1}^n \int_{t_{i-1}^n}^{t_i^n}  \l(\E| Z_u - Z_{t_{i-1}^n}|^2\r)^{1/2}\l( \E\l(\int_0^1 \p_x f(Z_{t_{i-1}^n} + v (Z_u - Z_{t_{i-1}^n}), \th)\,\df v\r)^2\r)^{1/2}\,\df u. 
\end{align*}
Since 
\[
(Z_t,Z_s)^\top \sim {\cal N}(\mb{0}, \Sigma(t,s)),\quad \Sigma(t,s) = 
\begin{pmatrix}
K(0) & K(t-s)\\ 
K(t-s) & K(0)
 \end{pmatrix}, 
\]
we see that
\begin{align*}
\E| Z_u - Z_{t_{i-1}^n}|^2 = 2\l[K(0)- K(u - t_{i-1}^n)\r]. 
\end{align*}
Then Corollary \ref{cor:ergodic} yields the convergence in probability for each $\th\in \ol{\Theta}$:
\begin{align}
\frac{1}{n} \sum_{i=1}^n f(Z_{t_{i-1}^n},\th)  \toP \E[f(Z_0,\th)] = \int_\R f(z,\th)\phi_{K(0)}(z)\,\df z. \label{lem:(i)}
\end{align}

As for the uniformity,  we shall confirm that
\begin{align}
\sup_{n\in \N} \E\l[\sup_{\th \in \ol{\Theta}}\l| \frac{1}{n}\sum_{i=1}^n \p_\th f(Z_{t_{i-1}^n},\th)\r|\r] < \infty,  \label{lem:(ii)}
\end{align}
which is easy to see by the condition \eqref{cond:ergod}. 

Then \eqref{lem:(i)} and \eqref{lem:(ii)} yield the consequence.
\end{proof}



\begin{lem}\label{lem:GZ}
Let $Z = (Z_t)_{t \ge 0} \sim GP(0,K)$ with A\ref{as:2} and  A\ref{as:3}, and let $G(x,y) :\R^2\to \R$ be a function such that $G(x,\cdot)\in C^2(\R)$ for each $x\in \R$.
Suppose that, for each $x\in \R$ and $k=1,2,\dots, l \ (l\ge 2)$, $\p_y^kG(x,y)$ is of polynomial growth w.r.t. $y$, $G(x,x) = 0$, and that  there exists a constant $M>0$ such that $|\p_y^{l+1}G(x,y)| \le M$ for all $x,y\in \R$. Then it holds that
\begin{align*}
\frac{1}{nh_n^2}\sum_{i=1}^n G(Z_{t_{i-1}^n}, Z_{t_i^n})
&\toP  \frac{\p^2 K(0)}{2 K(0)}  \int_\R \l[\p_y G(z,z)\,z - \p_y^2 G(z, z)K(0)\r] \, \phi_{K(0)}(z)\,\df z,
\end{align*}
under $h_n\to 0$ as $n\to \infty$. 
\end{lem}


\begin{proof}
Note that  $Z_{t_i} = Z_{t_{i-1}} + \Delta_i^n Z$
to apply Taylor's formula of $y \mapsto G(x,y)$ around $y = x$: 
\begin{align*}
G(Z_{t_{i-1}}, Z_{t_i}) 
&= G(Z_{t_{i-1}}, Z_{t_{i-1}} + \Delta_i^n Z) \\
&= \p_y G(Z_{t_{i-1}}, Z_{t_{i-1}}) \cdot \Delta_i^n Z 
+ \frac{1}{2} \p_y^2 G(Z_{t_{i-1}}, Z_{t_{i-1}}) \cdot (\Delta_i^n Z)^2 
+ R_i,
\end{align*}
where the remainder $R_i$ is given by
\[
R_i = \frac{1}{6} \p_y^3 G(Z_{t_{i-1}}, Z_{t_{i-1}} + \theta_i \Delta_i^n Z) \cdot (\Delta_i^n Z)^3
\quad \text{for some } \theta_i \in (0,1).
\]
Now we examine each term in the average
\[
\frac{1}{n h_n^2} \sum_{i=1}^n G(Z_{t_{i-1}}, Z_{t_i}) =: T_1 + T_2 + T_3,
\]
where $T_3 = \frac{1}{n h_n^2} \sum_{i=1}^n R_i$; 
\begin{align*}
T_1 &:= \frac{1}{n h_n^2} \sum_{i=1}^n \p_y G(Z_{t_{i-1}}, Z_{t_{i-1}}) \cdot \Delta_i^n Z; \quad 
T_2 := \frac{1}{2n h_n^2} \sum_{i=1}^n \p_y^2 G(Z_{t_{i-1}}, Z_{t_{i-1}}) \cdot (\Delta_i^n Z)^2. 
\end{align*}
We decompose $T_1$ as
\begin{align*}
T_1 = &= \frac{1}{n h_n^2} \sum_{i=1}^n \p_y G(Z_{t_{i-1}}, Z_{t_{i-1}}) \cdot \E[\Delta_i^n Z | Z_{t_{i-1}}] \\
&+ \frac{1}{n h_n^2} \sum_{i=1}^n \p_y G(Z_{t_{i-1}}, Z_{t_{i-1}}) \cdot \l( \Delta_i^n Z - \E[\Delta_i^n Z|Z_{t_{i-1}}] \r)\\
&=:A_n + B_n. 
\end{align*}
Since $\E[\Delta_i^n Z | Z_{t_{i-1}} = z] = \rho_n z$ with $\rho_n = \frac{K(h_n)}{K(0)} = 1 - \frac{1}{2} \p_t^2 K(0) h_n^2 + o(h_n^3)$, 
we have
\[
\E[\Delta_i^n Z|Z_{t_{i-1}}] = -\frac{1}{2} \p_t^2 K(0) \cdot \frac{h_n^2}{K(0)} Z_{t_{i-1}} + o(h_n^3).
\]
Therefore,
\[
A_n = -\frac{\p_t^2 K(0)}{2 K(0)} \cdot \frac{1}{n} \sum_{i=1}^n \p_y G(Z_{t_{i-1}}, Z_{t_{i-1}}) \cdot Z_{t_{i-1}} + o_P(1).
\]
Since the function $f(z) := \p_y G(z,z) \cdot z$ is of polynomial growth and satisfies the condition \eqref{cond:ergod}, Lemma~\ref{lem:ergod} implies that
\[
\frac{1}{n} \sum_{i=1}^n \p_y G(Z_{t_{i-1}}, Z_{t_{i-1}}) \cdot Z_{t_{i-1}} \toP \int_{\mathbb{R}} \p_y G(z,z) \cdot z \, \phi_{K(0)}(z) \,\df z.
\]
Hence,
\[
A_n \toP -\frac{\p_t^2 K(0)}{2 K(0)} \int_{\mathbb{R}} \p_y G(z,z) \cdot z \, \phi_{K(0)}(z) \,\df z.
\]
Next, we handle the centered term $B_n$. Note that 
\[
B_n= \frac{1}{n} \sum_{i=1}^n \p_y G(Z_{t_{i-1}}, Z_{t_{i-1}}) \cdot \wt{Z}_i^n, \quad \wt{Z}_i^n = \frac{\Delta_i^n Z - \E[\Delta_i^n Z| Z_{t_{i-1}}]}{h_n^2}, 
\]
where $\E[\wt{Z}_i^n] = 0$ and $\Var(\wt{Z}_i^n) =O(h_n^2)$. According to the Schwartz inequality, we have 
\[
B_n^2 \le \l(\frac{1}{n}\sum_{i=1}^n \p_y G(Z_{t_{i-1}}, Z_{t_{i-1}})^2 \r) \l(\frac{1}{n}\sum_{i=1}^n | \wt{Z}_i^n |^2  \r) = O_p(h_n^2)\toP0,\quad n\to \infty, 
\] 
by Lemma \ref{lem:ergod}. 



Thus, combining both parts, we conclude that 
\[
T_1 \toP -\frac{\p_t^2 K(0)}{2 K(0)} \int_{\mathbb{R}} \p_y G(z,z) \cdot z \, \phi_{K(0)}(z) \,\df z.
\]
As for $T_2$, the same argument leads us that 
\[
\frac{1}{2n h_n^2} \sum_{i=1}^n \p_y^2 G(Z_{t_{i-1}}, Z_{t_{i-1}}) \cdot (\Delta_i^n Z)^2
\toP \frac{\p^2 K(0)}{2} \int_{\mathbb{R}} \p_y^2 G(z,z) \, \phi_{K(0)}(z) \,\df z, 
\]
and the details are omitted. 

As for $T_3$, since $|\p_y^3 G(x,y)| \le M$ and $\E[|\Delta_i^n Z|^3] = O(h_n^3)$, it follows from Lemma \ref{lem:avg-conv} that 
\[
\l|\frac{1}{n h_n^2} \sum_{i=1}^n R_i \r| \le M \frac{h_n}{n}\sum_{i=1}^n \l(\frac{\D_i^n Z}{h_n}\r)^3=O_p(h_n) \to 0.
\]
As a result, we obtain that 
\[
\frac{1}{n h_n^2} \sum_{i=1}^n G(Z_{t_{i-1}}, Z_{t_i}) 
\toP \frac{\p^2 K(0)}{2 K(0)} \int_{\mathbb{R}} \l[ \p_y G(z,z) \cdot z - \p_y^2 G(z,z) \cdot K(0) \r] \phi_{K(0)}(z) \,\df z.
\]
\end{proof}



\section{Auxiliary Lemmas}\label{app:B}
In this section, we assume Conditions A\ref{as:2}--A\ref{as:4} without further mention.

\begin{lem}\label{lem:f-DX}
Let $f:\mathbb{R}\times \ol{\Theta} \to \mathbb{R}$, be continuous and of polynomial growth uniformly in $\th \in \ol{\Theta}$:
\begin{align}
|\p^k_x\p_\th^lf(x,\th)| \le C(1 + |x|^p), \quad  x \in \mathbb{R},  \label{cond:f-DX}
\end{align}
for integers $k$ and $l$ with $0\le k+l\le 1$, $C>0$ and $p \ge 1$, and suppose the assumption B\ref{as:mu} in the model \eqref{model}. 
Then, it holds that 
\begin{align*}
\frac{1}{n} \sum_{i=1}^n f \l(\frac{\D_i^n X - \mu_\xi(t_{i-1}) h_n}{h_n},\th\r) \toP \int_\R f(z,\th) \phi_{-\p^2_t K(0)}(z)\,\df z, 
\end{align*}
as $n\to \infty$, uniformly in $\th \in \ol{\Theta}$. 
\end{lem}

\begin{proof}
First, we fix $\th\in \Theta$, and define the scaled increments:
\[
\widetilde{Y}_i^n(\xi) := \frac{\Delta_i^n X - \mu_\xi(t_{i-1})h_n }{h_n}, \qquad Y_i^n := \frac{\Delta_i^n Z}{h_n}.
\]
We have the decomposition
\begin{align*}
\frac{1}{n} \sum_{i=1}^n f(\widetilde{Y}_i^n(\xi),\th) = \frac{1}{n} \sum_{i=1}^n f(Y_{i-1}^n,\th) + \frac{1}{n} \sum_{i=1}^n \l( f(\widetilde{Y}_i^n(\xi),\th) - f(Y_{i-1}^n,\th) \r).
\end{align*}
We first prove that the second term converges to zero in probability uniformly in $\theta$.

By the mean value theorem, there exists $y_i^n$ between $\widetilde{Y}_i^n(\xi)$ and $Y_{i-1}^n$ such that
\[
|f(\widetilde{Y}_i^n(\xi),\theta) - f(Y_{i-1}^n,\theta)| = |\partial_x f(y_i^n,\theta)||\widetilde{Y}_i^n(\xi) - Y_{i-1}^n|. 
\]
Using the polynomial growth condition, we have
\[
|\partial_x f(y_i^n, \th)| \le C(1 + |y_i^n|^p) \quad \text{for some constant } C>0.
\]
Next, observe that
\begin{align*}
\widetilde{Y}_i^n(\xi) - Y_{i-1}^n &= \frac{1}{h_n} \l( \Delta_i^n X - \mu_\xi(t_{i-1}) h_n - \Delta_i^n Z \r) \\
&= \frac{1}{h_n} \l( \int_{t_{i-1}}^{t_i} \mu_\xi(s)\,\df s - \mu_\xi(t_{i-1}) h_n \r) \\
&= \frac{1}{h_n} \int_{t_{i-1}}^{t_i} (\mu_\xi(s) - \mu_\xi(t_{i-1}))\,\df s.
\end{align*}
By the continuity of $\mu_\xi$ and the mean value theorem, there exists $\eta_i^n \in [t_{i-1}, t_i]$ such that
\[
\mu_\xi(s) - \mu_\xi(t_{i-1}) = \partial_t \mu_\xi(\eta_i^n)(s - t_{i-1}),
\]
thus, by the condition B\ref{as:mu}, 
\begin{align*}
|\widetilde{Y}_i^n(\xi) - Y_{i-1}^n| &\le \frac{1}{h_n} \int_{t_{i-1}}^{t_i} |\partial_t \mu_\xi(\eta_i^n)||s - t_{i-1}|\,\df s \\
&\le \frac{\sup_{t>0,\th\in \ol{\Theta}}|\partial_t \mu_\xi(t)|}{h_n} \int_0^{h_n} u\,\df u  = O_p(h_n). 
\end{align*}
Hence, $|\widetilde{Y}_i^n(\xi) - Y_{i-1}^n| = O_p(h_n)$ uniformly in $i$.

Moreover, since $y_i^n$ lies between $Y_{i-1}^n$ and $\widetilde{Y}_i^n(\xi)$, which are Gaussian variables, it follows that $|y_i^n|$ also has a Gaussian tail. 
This implies that 
\[
\E\l[ (1 + |y_i^n|^p) |\widetilde{Y}_i^n(\xi) - Y_{i-1}^n| \r] \le C h_n, 
\]
uniformly in $i$ and $\theta$.
Therefore,
\[
\frac{1}{n} \sum_{i=1}^n |f_\theta(\widetilde{Y}_i^n(\xi)) - f_\theta(Y_{i-1}^n)| \toP 0, 
\]
for each $\th\in \Theta$. Moreover, the tightness of $\frac{1}{n} \sum_{i=1}^n |f_\theta(\widetilde{Y}_i^n(\xi)) - f_\theta(Y_{i-1}^n)| $ is also easy to see from the condition \eqref{cond:f-DX}, e.g., by the same argument as in the proof of Lemma \ref{lem:avg-conv}. 
Hence the above convergence is indeed uniform in $\theta \in \overline{\Theta}$. 
As a consequence, Lemma \ref{lem:avg-conv} completes the proof. 
\end{proof}





\begin{lem}\label{lem:score-vector}
Suppose that Assumptions B\ref{as:DRI-mu}, B\ref{as:DRI-mu2} and B\ref{as:DRI-mu3} hold. 
Then, it follows for the block diagonals $D_n := \mathrm{diag}\l( h_n^{-1/2} I_p,  \sqrt{n}I_q\r)$  and $C_n := \mathrm{diag}\l( nh_n I_p, I_q\r)$ that 
\[
C_nD_n \partial_\theta \ell_n(\theta_0)
=
\begin{pmatrix}
n\sqrt{h_n}\partial_\xi \ell_n(\theta_0) \\
\sqrt{n} \partial_\sigma \ell_n(\theta_0)
\end{pmatrix}
\toD \mathcal{N}(0, J(\theta_0)),
\]
where  
\[
J(\theta_0) := 
\begin{pmatrix}
 -\frac{2}{\partial_t^2 K(0)}\int_0^\infty \{\p_\xi \mu_{\xi}(t)\}^{\otimes 2}\, \df t &0 \\
0 & \l( \frac{1}{2} \p_\s \log (-\partial_t^2 K_{\sigma_0}(0) \r)^{\otimes 2} 
\end{pmatrix}. 
\]
\end{lem}


\begin{proof}
From the contrast function:
\[
\ell_n(\xi, \sigma) := \frac{1}{n} \sum_{i=1}^n \l\{
\frac{(\Delta_i^n X - h_n \mu_\xi(t_{i-1}))^2}{2 [K_\sigma(0) - K_\sigma(h_n)]} +
\log \l( 2 h_n^{-2} [K_\sigma(0) - K_\sigma(h_n)] \r)
\r\}, 
\]
the score vector is given by
\begin{align*}
\partial_\xi \ell_n(\xi, \sigma)
&= - \frac{h_n}{n v_n(\s) } \sum_{i=1}^n (\Delta_i^n X - h_n \mu_\xi(t_{i-1})) \cdot \partial_\xi \mu_\xi(t_{i-1}), \\
\partial_\sigma \ell_n(\xi, \sigma)
&= \sum_{i=1}^n \l( -\frac{(\Delta_i^n X - h_n \mu_\xi(t_{i-1}))^2}{2 v_n^2(\s)} + \frac{1}{v_n(\s) } \r) \cdot \frac{1}{n} \partial_\sigma v_n(\s) .
\end{align*}
where $v_n(\s) := K_\sigma(0) - K_\sigma(h_n)$.
Noticing that $\D_i^n X = \D_i^n Z + \delta_i^n$ with
\[
\delta_i^n := \int_{t_{i-1}}^{t_i} \mu_{\xi_0}(s)\, \mathrm{d}s - \mu_{\xi_0}(t_{i-1}) h_n 
= \frac{1}{2} \p_t \mu_{\xi_0}(t_{i-1}) h_n^2 + o(h_n^2) = O(h_n^2), 
\]
and 
\[
\frac{h_n}{v_n(\sigma_0)} = \frac{2}{h_n \partial_t^2 K(0)} + o(1).
\]
since $v_n(\sigma_0) = \frac{1}{2} h_n^2 \partial_t^2 K(0) + o(h_n^3)$ by A\ref{as:3}, we obtain that 
\begin{align*}
\p_\xi \ell_n(\theta_0)
&= -\l(\frac{2}{h_n \p_t^2K(0)} + o(1)\r)\Bigg[\frac{1}{n}\sum_{i=1}^n \D_i^nZ \p_\xi\mu_{\xi_0}(t_{i-1}) - \frac{1}{n}\sum_{i=1}^n\p_\xi\mu_{\xi_0}(t_{i-1}) \cdot O(h_n^2) \\
&\qquad + \frac{1}{n}\sum_{i=1}^n \mu_{\xi_0}(t_{i-1}) \p_\xi\mu_{\xi_0}(t_{i-1})h_n \Bigg] \\
&= -\frac{2}{nh_n \p_t^2K(0)} \sum_{i=1}^n \D_i^nZ \p_\xi\mu_{\xi_0}(t_{i-1}) + O_p\l(\frac{1}{n}\r). 
\end{align*}
 under the DRI condition B\ref{as:DRI-mu} and B\ref{as:DRI-mu2}.
Now multiply both sides by $n \sqrt{h_n} $:
\begin{align*}
n \sqrt{h_n} \cdot \p_\xi \ell_n(\theta_0)
&= -\frac{2}{\p_t^2K(0)} \sum_{i=1}^n \frac{1}{\sqrt{h_n}}\D_i^nZ \p_\xi\mu_{\xi_0}(t_{i-1}) + O_p\l(\sqrt{h_n}\r).
\end{align*}
Next, it follows that 
\begin{align*}
\sqrt{n}  \partial_\sigma \ell_n(\theta_0) &= \frac{\partial_\sigma v_n(\s_0)}{\sqrt{n} } \sum_{i=1}^n \l( \frac{1}{v_n(\s_0)} - \frac{1}{2 v^2_n(\s_0)} (\Delta_i^n Z)^2 \r).
\end{align*}
using $\partial_\sigma v_n(\sigma_0) = - \frac{h_n^2}{2} \partial_t \partial_\sigma K(0) + o(h_n^3)$ and the above expansion of $v_n$, we find:
\[
\frac{\partial_\sigma v_n(\sigma_0)}{v_n(\sigma_0)} = \frac{\partial_\sigma \partial_t^2 K_{\s_0}(0)}{\partial_t^2 K(0)} + o(h_n).
\]
Hence,
\[
\partial_\sigma \ell_n(\theta_0) = \l( \frac{\partial_\sigma \partial_t^2K_{\s_0}(0)}{\partial_t^2 K(0)} + o(h_n) \r) \cdot \frac{1}{n} \sum_{i=1}^n \l(1 - \frac{(\Delta_i^n Z)^2}{2 v_n(\sigma_0)}\r).
\]
Since each summand is $O_p(1)$ by Lemma \ref{lem:avg-conv}, the remainder term becomes $o_p(h_n)$. 

As a summary, letting 
\[
X_i^n := \sqrt{h_n} \Delta_i^n Z \cdot \partial_\xi \mu_{\xi_0}(t_{i-1}),
\quad
Y_i^n := \frac{1}{\sqrt{n}} \l( 1 - \frac{(\Delta_i^n Z)^2}{2 v_n(\s_0)} \r), 
\]
we have that 
\begin{align*}
n \sqrt{h_n}  \partial_\xi \ell_n(\theta_0) &= - \frac{2}{\partial_t^2 K(0)} \sum_{i=1}^n X_i^n + o_p(1), \\
\sqrt{n} \partial_\sigma \ell_n(\theta_0) &= \frac{\partial_\sigma \partial_t^2 K_{\s_0}(0)}{\partial_t^2 K(0)} \sum_{i=1}^n Y_i^n + o_p(1).
\end{align*}

We now verify the conditions of Theorem~2.1 in Neumann~\cite{n13} for the triangular arrays $\{X_i^n\}$ and $\{Y_i^n\}$ defined by:
\[
X_i^n := \frac{1}{\sqrt{h_n}}\Delta_i^n Z \cdot \partial_\xi \mu_{\xi_0}(t_{i-1}),
\quad
Y_i^n := \frac{1}{\sqrt{n}} \l( 1 - \frac{(\Delta_i^n Z)^2}{2 v_n(\sigma_0)} \r).
\]

Now, notice that 
\begin{align*}
D_n \partial_\theta \ell_n(\theta_0) 
&=\mathrm{diag}\l(- \frac{2}{\partial_t^2 K(0)}I_p,  \frac{1}{\partial_t^2 K(0)}I_q\r) 
\sum_{i=1}^n \begin{pmatrix} X_i^n \\Y_i^n\partial_\sigma \partial_t^2 K_{\s_0}(0)\end{pmatrix}  + o_p(1), 
\end{align*}
To show the weak convergence of $S_n$, we apply the Cram\'er--Wold device. 
That is, for any $(a,b)^\top \in \mathbb{R}^{p+q}$ with $a \in \mathbb{R}^p$, $b \in \mathbb{R}^q$, we denote by 
\[
S_n := \sum_{i=1}^n \l( a^\top X_i^n + b^\top Y_i^n\partial_\sigma \partial_t^2 K_{\s_0}(0) \r).
\]
and show the weak convergence of $S_n$ by applying Thereom 2.1 in Neuman \cite{n13}, the CLT for trriangular arrays. 

\paragraph{(i) Mean zero.}
Since $Z$ is a centered Gaussian process, $\E[\Delta_i^n Z] = 0$ and $\E[(\Delta_i^n Z)^2] = v_n(\sigma_0)$, we have 
$\E[X_i^n] = 0$ and $\E[Y_i^n] = 0$, and $\E[S_n]=0$. 

\paragraph{(ii) Variance convergence.}
For the $X$-component:
\begin{align*}
\sum_{i=1}^n \E[(X_i^n)^{\otimes 2}] &= h_n^{-1}\sum_{i=1}^n \E[(\Delta_i^n Z)^2] \cdot \l\{ \partial_\xi \mu_{\xi_0}(t_{i-1}) \r\}^{\otimes 2} \\
&= v_n(\sigma_0) h_n^{-1}\sum_{i=1}^n \l\{ \partial_\xi \mu_{\xi_0}(t_{i-1}) \r\}^{\otimes 2} \to \frac{1}{2} \partial_t^2 K(0) \int_0^\infty \l\{ \partial_\xi \mu_{\xi_0}(t) \r\}^{\otimes 2} dt.
\end{align*}
by the condition B\ref{as:DRI-mu2}. 
For the $Y$-component:
\begin{align*}
\sum_{i=1}^n \E[(Y_i^n)^2] &= \frac{1}{n} \sum_{i=1}^n \operatorname{Var}\l( 1 - \frac{(\Delta_i^n Z)^2}{2 v_n(\sigma_0)} \r) \\
&= \frac{1}{n} \sum_{i=1}^n \frac{1}{4} \operatorname{Var}\l( \frac{(\Delta_i^n Z)^2}{v_n(\sigma_0)} \r) = \frac{1}{4} \cdot \frac{1}{n} \cdot 2n = \frac{1}{2}.
\end{align*}
Moreover, 
\[
\E[X_i^n Y_i^n] = \frac{1}{\sqrt{n h_n}} \cdot \E\l[ \Delta_i^n Z \cdot \l( 1 - \frac{(\Delta_i^n Z)^2}{2 v_n(\sigma_0)} \r) \r] \cdot \partial_\xi \mu_{\xi_0}(t_{i-1}).
\]
Since $\E[\Delta_i^n Z] = \E[\Delta_i^n Z^3] =0$, we obtain that $\E[X_i^n Y_i^n] = 0$.
As a result, we see that 
\[
\E[(S_n)^2] \to a^\top \l(\frac{2}{\partial_t^2 K(0)}  \int_0^\infty \l\{ \partial_\xi \mu_{\xi_0}(t) \r\}^{\otimes 2}\,\df t\r) a + b^\top \frac{1}{2}\l(\partial_\sigma \partial_t^2 K_{\s_0}(0)\r)^{\otimes 2} b. 
\]

\paragraph{(i) Lyapnov condition.}
Since $X_i^n$ and $Y_i^n$ are centered Gaussian (or polynomial transformations thereof), we can compute their fourth moments explicitly:
\[
\E[|X_i^n|^4] = \l( \frac{1}{h_n} \r)^2 \cdot 3 v_n^2 \cdot \|\partial_\xi \mu(t_{i-1})\|^4 = O\l( h_n^2\r),
\]
\[
\E[|Y_i^n|^4] = \frac{1}{n^2} \cdot \E\l[ \l(1 - \frac{1}{2} \chi^2(1) \r)^4 \r] = O\l( \frac{1}{n^2} \r).
\]
Therefore,
\[
\sum_{i=1}^n \E[|S_{n,i}|^4] \le 8 \sum_{i=1}^n \l( \E[|a^\top X_i^n|^4] + \E[|b^\top Y_i^n|^4] \r) = O\l(h_n^2 +  \frac{1}{n^2} \r) \to 0.
\]
\paragraph{(ii) Mixing covariance bounds.}
We apply Lemmas~\ref{lem:cov-bound-summable} and~\ref{lem:cov-bound-summable-Y} below, which establish uniform bounds for the covariance of nonlinear functionals of the triangular arrays $\{X_i^n\}$ and $\{Y_i^n\}$ under the assumptions A\ref{as:2} and A\ref{as:3}. In particular, for any bounded measurable function $g$ with $\|g\|_\infty \le 1$, and indices $1 \le s_1 < \dots < s_u < t_1 \le n$, the inequalities
\[
\l| \Cov\l( g(X_{s_1}^n, \dots, X_{s_u}^n) X_{s_u}^n,\; X_{t_1}^n \r) \r|
\le \l( \E|X_{s_u}^n|^2 + \E|X_{t_1}^n|^2 + \frac{1}{n} \r) \cdot \theta_{r},
\]
\[
\l| \Cov\l( g(Y_{s_1}^n, \dots, Y_{s_u}^n),\; Y_{t_1}^n Y_{t_2}^n \r) \r|
\le \l( \E|Y_{t_1}^n|^2 + \E|Y_{t_2}^n|^2 + \frac{1}{n} \r) \cdot \theta_{r},
\quad \text{with } r = t_1 - s_u,
\]
hold for a summable sequence $\{\theta_r\} \in \ell^1$ depending on the covariance kernel $K$. This verifies condition (2.6) of Neumann~\cite{n13}. The summability of $\{\theta_r\}$ follows from the integrability condition $K \in L^1(\R)$ imposed in A\ref{as:3}, ensuring that the dependence decays sufficiently fast to guarantee asymptotic independence in the triangular arrays.

\end{proof}

\begin{remark}
The sequence $\{\theta_r\}_{r \in \mathbb{N}}$ defined in Lemma~\ref{lem:score-vector} and the lemmas below depends on $n$ through the mesh size $h_n$.
This dependence is admissible in the framework of Neumann \cite{n13}, as the central limit theorem for triangular arrays of weakly dependent variables requires only that the dependence coefficients (such as $\theta_r$) satisfy a uniform summability condition over $n$:
\[
\sup_{n \in \mathbb{N}} \sum_{r=1}^\infty \theta_r^{(n)} < \infty.
\]
In our case, this is ensured by the integrability condition $K \in L^1((0,\infty))$.
\end{remark}

\begin{lem}
\label{lem:cov-bound-summable}
Let $\{X_i^n\}_{1 \le i \le n}$ be the triangular array defined by
\[
X_i^n := \frac{1}{\sqrt{h_n}} \Delta_i^n Z \cdot \partial_\xi \mu_{\xi_0}(t_{i-1}),
\]
where $\Delta_i^n Z := Z_{t_i} - Z_{t_{i-1}}$ and define for $r \in \N$ the sequence
\[
\theta_r := \frac{ \displaystyle \int_0^{h_n} \!\! \int_0^{h_n} |K(r h_n + u - v)|\,\df u\,\df v }{
\displaystyle 2 \int_0^{h_n} \!\! \int_0^{h_n} |K(u - v)|\,\df u\,\df v + h_n }.
\]
Then, for any bounded measurable function $g$ with $\|g\|_\infty \le 1$, 
and any indices $1 \le s_1 < \dots < s_u < t_1 \le n$ with $r := t_1 - s_u$, it holds under B\ref{as:DRI-mu3} that
\[
\l| \Cov\l( g(X_{s_1}^n, \dots, X_{s_u}^n) X_{s_u}^n,\; X_{t_1}^n \r) \r|
\le \l( \E|X_{s_u}^n|^2 + \E|X_{t_1}^n|^2 + \frac{1}{n} \r) \cdot \theta_r,  
\]
and 
\[
\l| \Cov\l( g(X_{s_1}^n, \dots, X_{s_u}^n),\; X_{t_1}^n X_{t_2}^n \r) \r|
\le \l( \E|X_{t_1}^n|^2 + \E|X_{t_2}^n|^2 + \frac{1}{n} \r) \cdot \theta_r,  
\]
Moreover, $\sum_{r=1}^\infty \theta_r < \infty$.
\end{lem}

\begin{proof}
First, note that
\[
X_i^n = \frac{1}{\sqrt{h_n}} \cdot \partial_\xi \mu_{\xi_0}(t_{i-1}) \cdot \Delta_i^n Z,
\]
so that
\[
\Cov\l( g(\cdots) X_{s_u}^n,\; X_{t_1}^n \r)
= \frac{1}{h_n} \cdot \partial_\xi \mu_{\xi_0}(t_{s_u - 1}) \cdot \partial_\xi \mu_{\xi_0}(t_{t_1 - 1}) \cdot \E\l[ g(\cdots) \Delta_{s_u}^n Z \cdot \Delta_{t_1}^n Z \r].
\]
Since $|g| \le 1$, it follows from the Cauchy-Schwarz inequality that
\[
\l| \E\l[ g(\cdots) \Delta_{s_u}^n Z \cdot \Delta_{t_1}^n Z \r] \r|
\le \E\l| \Delta_{s_u}^n Z \cdot \Delta_{t_1}^n Z \r|
\le \int_0^{h_n} \int_0^{h_n} |K(r h_n + u - v)|\,\df u\,\df v.
\]
Similarly, the variance terms can be bounded as
\[
\E|X_i^n|^2 \le \frac{C_\mu^2}{h_n} \cdot \int_0^{h_n} \int_0^{h_n} |K(u - v)|\,\df u\,\df v, 
\]
where $C_\mu$ is a constant such that $|\g(t)|\le C_\mu$ in the assumption B\ref{as:DRI-mu3}. 
Therefore, the desired inequality holds with $\theta_r$ as defined.

To prove the summability of $\{\theta_r\}$, we use the change of variables $x = r h_n + s$ with $s \in [-h_n, h_n]$, and estimate
\[
\sum_{r=1}^\infty \theta_r
\le C \sum_{r=1}^\infty \int_0^{h_n} \int_0^{h_n} |K(r h_n + u - v)|\,\df u\,\df v
\le C' \int_0^\infty |K(x)|\,\df x < \infty,
\]
where $C, C'$ are constants independent of $r$. 

Moreover, since 
\begin{align*}
\l| \Cov\l( g(\cdots),\; X_{t_1}^n X_{t_2}^n \r) \r|
&\le \E\l| g(\cdots) \cdot X_{t_1}^n X_{t_2}^n \r| \le \E\l| X_{t_1}^n X_{t_2}^n \r| \\
&= \frac{1}{n h_n} \cdot |\partial_\xi \mu_{\xi_0}(t_{t_1 - 1})| \cdot |\partial_\xi \mu_{\xi_0}(t_{t_2 - 1})| \cdot \E\l| \Delta_{t_1}^n Z \cdot \Delta_{t_2}^n Z \r|.
\end{align*}
the similar argument leads us to the consequence. 
\end{proof}


\begin{lem}
\label{lem:cov-bound-summable-Y}
Let $\{Y_i^n\}_{1 \le i \le n}$ be the triangular array defined by
\[
Y_i^n := \frac{1}{\sqrt{n}} \l( 1 - \frac{(\Delta_i^n Z)^2}{2 v_n(\sigma_0)} \r),
\]
Define for $r \in \N$ the sequence
\[
\theta_r := \frac{ \displaystyle \int_0^{h_n} \!\! \int_0^{h_n} |K(r h_n + u - v)|\,\df u\,\df v }{
\displaystyle 2 \int_0^{h_n} \!\! \int_0^{h_n} |K(u - v)|\,\df u\,\df v + h_n }.
\]
Then, for any bounded measurable function $g$ with $\|g\|_\infty \le 1$, and any indices $1 \le s_1 < \dots < s_u < t_1 \le n$ with $r := t_1 - s_u$, it holds that
\[
\l| \Cov\l( g(Y_{s_1}^n, \dots, Y_{s_u}^n) Y_{s_u}^n,\; Y_{t_1}^n \r) \r|
\le \l( \E|Y_{s_u}^n|^2 + \E|Y_{t_1}^n|^2 + \frac{1}{n} \r) \cdot \theta_r,
\]
and
\[
\l| \Cov\l( g(Y_{s_1}^n, \dots, Y_{s_u}^n),\; Y_{t_1}^n Y_{t_2}^n \r) \r|
\le \l( \E|Y_{t_1}^n|^2 + \E|Y_{t_2}^n|^2 + \frac{1}{n} \r) \cdot \theta_r.
\]
Moreover, $\sum_{r=1}^\infty \theta_r < \infty$.
\end{lem}

\begin{proof}
We first observe that $Y_i^n$ can be written as
\[
Y_i^n = \frac{1}{\sqrt{n}} \l( 1 - \frac{(\Delta_i^n Z)^2}{2 v_n(\sigma_0)} \r),
\]
so that
\[
\E|Y_i^n|^2 = \frac{1}{n} \cdot \Var\l( \frac{(\Delta_i^n Z)^2}{2 v_n(\sigma_0)} \r)
\le \frac{C}{n} \cdot \Var\l( (\Delta_i^n Z)^2 \r),
\]
where the last variance is controlled by the fourth moment of $\Delta_i^n Z$. Since $Z$ is a centered stationary Gaussian process, we obtain
\[
\E[(\Delta_i^n Z)^4] \le C \cdot \l( \int_0^{h_n} \!\! \int_0^{h_n} |K(u - v)|\,\df u\,\df v \r)^2,
\]
and hence
\[
\E|Y_i^n|^2 \le \frac{C'}{n} \cdot \int_0^{h_n} \!\! \int_0^{h_n} |K(u - v)|\,\df u\,\df v.
\]
This upper bound motivates the choice of denominator in $\theta_r$ so that the quantity $\E|Y_i^n|^2$ is uniformly absorbed.
Now, for $\|g\|_\infty \le 1$, we apply the Cauchy--Schwarz inequality:
\[
\l| \Cov\l( g(\cdots) Y_{s_u}^n,\; Y_{t_1}^n \r) \r|
\le \E\l| Y_{s_u}^n Y_{t_1}^n \r|.
\]
Using the above representation, we have
\[
\E\l| Y_{s_u}^n Y_{t_1}^n \r|
= \frac{1}{n} \cdot \E\l| \l( 1 - \frac{(\Delta_{s_u}^n Z)^2}{2 v_n(\sigma_0)} \r)
\l( 1 - \frac{(\Delta_{t_1}^n Z)^2}{2 v_n(\sigma_0)} \r) \r|.
\]
The dominant term arises from the covariance between $(\Delta_{s_u}^n Z)^2$ and $(\Delta_{t_1}^n Z)^2$, which is controlled by
\[
\E\l| (\Delta_{s_u}^n Z)^2 \cdot (\Delta_{t_1}^n Z)^2 \r|
\le C \cdot \int_0^{h_n} \!\! \int_0^{h_n} |K(r h_n + u - v)|\,\df u\,\df v.
\]
Therefore, the bound with $\theta_r$ holds.
The same reasoning applies to
\[
\l| \Cov\l( g(\cdots),\; Y_{t_1}^n Y_{t_2}^n \r) \r|
\le \E\l| Y_{t_1}^n Y_{t_2}^n \r|,
\]
which is again controlled by the same integral involving $K(r h_n + u - v)$. Finally, since
\[
\sum_{r=1}^\infty \theta_r \le C \int_0^\infty |K(x)|\,\df x < \infty. 
\]
Hence, the lemma is proved.
\end{proof}


\begin{lem}\label{lem:fY}
Let $f(x,\th): \R\times \ol{\Theta} \to \R$ be a measureble function such that 
\begin{align}
\sup_{\th \in \ol{\Theta}}|\p_x^k \p_\th^l f(x,\th)| \lesssim 1 + |x|^C, \label{cond:ergod2}
\end{align}
for integers $k$ and $l$ with $0\le k+l\le 1$. 
Then, it holds under the assumption B\ref{as:DRI-mu3} that
\[
\sup_{\th \in \ol{\Theta}}\l|\frac{1}{n} \sum_{i=1}^n f(Y_{i-1}^n,\th) - \int_\R f(z,\th)\phi_{K(0)}(z)\,\df z\r|\toP 0,
\]
under $h_n \to 0$ and $nh_n\to \infty$ as $n\to \infty$, where $Y_i^n= X_{t_i} - \int_0^{t_i} \mu_{\wh{\xi}_n}(s)\,\df s$ and 
$\wh{\xi}_n$ is a consistent estimator for $\xi_0$: $\wh{\xi}_n\toP \xi_0$. 
\end{lem}

\begin{proof}
By the mean value theorem and B\ref{as:DRI-mu3}, we have
\[
\l| \mu_{\xi_0}(s) - \mu_{\widehat{\xi}_n}(s) \r|
\le \sup_{\xi \in \Xi} |\partial_\xi \mu_\xi(s)| \cdot |\widehat{\xi}_n - \xi_0|
\le \g(s) \cdot |\widehat{\xi}_n - \xi_0|,
\]
for each $s \ge 0$. Hence,
\[
\sup_{i \le n} |Y_{i-1}^n - Z_{t_{i-1}}|
\le |\widehat{\xi}_n - \xi_0| \cdot \int_0^{t_n} \g(s)\,\df s
\le |\widehat{\xi}_n - \xi_0| \cdot \|\g\|_{L^1([0,\infty))} \toP 0.
\]

Now, using the growth condition~\eqref{cond:ergod2} and the mean value theorem again, we obtain
\[
|f(Y_{i-1}^n,\theta) - f(Z_{t_{i-1}},\theta)|
\lesssim |Y_{i-1}^n - Z_{t_{i-1}}| \cdot \l(1 + |Z_{t_{i-1}}|^C + |Y_{i-1}^n|^C\r).
\]
Since $Z_{t_{i-1}}$ has finite moments and $Y_{i-1}^n \toP Z_{t_{i-1}}$ uniformly in $i$, it follows by dominated convergence that
\[
\sup_{\theta \in \overline{\Theta}} \l| \frac{1}{n} \sum_{i=1}^n f(Y_{i-1}^n,\theta) - \frac{1}{n} \sum_{i=1}^n f(Z_{t_{i-1}},\theta) \r| \toP 0.
\]

Finally, Lemma~\ref{lem:ergod} implies
\[
\sup_{\theta \in \overline{\Theta}} \l| \frac{1}{n} \sum_{i=1}^n f(Z_{t_{i-1}},\theta) - \int f(z,\theta)\,\phi_{K(0)}(z)\,\df z \r| \toP 0,
\]
and the claim follows by the triangle inequality.
\end{proof}


\begin{lem}\label{lem:GY}
Let $G(x,y,\th) :\R^2\times \Theta\to \R$ be a function such that $G(x,\cdot,\th)\in C^2(\R)$ for each $x\in \R$ and $\th\in \Theta$.
Suppose that, for each $x\in \R$ and $k=1,2,\dots, l \ (l\ge 2)$, $\p_y^kG(x,y,\th)$ is of polynomial growth w.r.t. $y$ uniformly in $\th\in \Theta$, $G(x,x,\th) = 0$, and that  there exists a constant $M>0$ such that $|\p_y^{l+1}G(x,y,\th)| \le M$ for all $x,y\in \R$ and $\th \in \Theta$. Then it holds that 
\begin{align*}
\frac{1}{nh_n^2}\sum_{i=1}^n G(Y^n_{i-1}, Y_i^n,\th)
&\toP  \frac{\p^2 K(0)}{2 K(0)}  \int_\R \l[\p_y G(z,z)\,z - \p_y^2 G(z, z)K(0)\r] \, \phi_{K(0)}(z)\,\df z,
\end{align*}
uniformly in $\th \in \ol{\Theta}$ under $h_n\to 0$ as $n\to \infty$, where $Y_i^n= X_{t_i} - \int_0^{t_i} \mu_{\wh{\xi}_n}(s)\,\df s$ and 
$\wh{\xi}_n$ is a consistent estimator for $\xi_0$: $\wh{\xi}_n\toP \xi_0$. 
\end{lem}
  
\begin{proof}
Let us fix \( \theta = (\xi, \sigma) \in \ol{\Theta} \). Define
\[
T_n(\theta) := \frac{1}{n h_n^2} \sum_{i=1}^n G\l(Y_{i-1}^n, Y_i^n, \theta\r),
\]
where \( Y_i^n := X_{t_i} - \int_0^{t_i} \mu_{\wh{\xi}_n}(s)\,\df s \) with a consistent estimator \( \wh{\xi}_n \toP \xi_0 \).

By the model definition \( X_t = Z_t + \int_0^t \mu_{\xi_0}(s)\,\df s \), we have:
\[
X_{t_i} - X_{t_{i-1}} = Z_{t_i} - Z_{t_{i-1}} + \int_{t_{i-1}}^{t_i} \mu_{\xi_0}(s)\,\df s.
\]
Thus,
\[
Y_i^n - Y_{i-1}^n = (Z_{t_i} - Z_{t_{i-1}}) + r_i^n, \quad\text{where } r_i^n := \int_{t_{i-1}}^{t_i} (\mu_{\xi_0}(s) - \mu_{\wh{\xi}_n}(s))\,\df s.
\]
Using Taylor expansion of \( y \mapsto G(x,y,\theta) \) around \( y = x \), we get
\begin{align*}
G(x,y,\theta) &= \p_y G(x,x,\theta)(y - x) + \frac{1}{2} \p_y^2 G(x,x,\theta)(y - x)^2 + R(x,y,\theta),
\end{align*}
with the remainder \( R(x,y,\theta) = \frac{1}{6} \p_y^3 G(x, \zeta, \theta)(y - x)^3 \) for some \( \zeta \) between \( x \) and \( y \). Then,
\begin{align*}
T_n(\theta) &= \frac{1}{n h_n^2} \sum_{i=1}^n \Bigg\{ \p_y G(Y_{i-1}^n, Y_{i-1}^n, \theta)(Y_i^n - Y_{i-1}^n) + \frac{1}{2} \p_y^2 G(Y_{i-1}^n, Y_{i-1}^n, \theta)(Y_i^n - Y_{i-1}^n)^2 + R_i^n(\theta) \Bigg\} \\
&=: A_n(\theta)  + B_n(\theta) + C_n(\theta). 
\end{align*}
with \( R_i^n(\theta) \) as above. We now evaluate \( A_n(\theta) \) in detail:
\begin{align*}
A_n(\theta) &= \frac{1}{n h_n^2} \sum_{i=1}^n \p_y G(Y_{i-1}^n, Y_{i-1}^n, \theta)(\Delta_i^n Z + r_i^n) \\
&= \frac{1}{n h_n^2} \sum_{i=1}^n \p_y G(Z_{t_{i-1}}, Z_{t_{i-1}}, \theta) \Delta_i^n Z + R'_n(\theta),
\end{align*}
where the error term \( R'_n(\theta) \) is decomposed as
\begin{align*}
R'_n(\theta) &= \frac{1}{n h_n^2} \sum_{i=1}^n \l[ \p_y G(Y_{i-1}^n, Y_{i-1}^n, \theta) - \p_y G(Z_{t_{i-1}}, Z_{t_{i-1}}, \theta) \r] \Delta_i^n Z \\
&\quad + \frac{1}{n h_n^2} \sum_{i=1}^n \p_y G(Y_{i-1}^n, Y_{i-1}^n, \theta) r_i^n.
\end{align*}
For the first term, using the integral form of the mean value theorem, we write:
\begin{align*}
\p_y G(Y_{i-1}^n, Y_{i-1}^n, \theta) - \p_y G(Z_{t_{i-1}}, Z_{t_{i-1}}, \theta) = \int_0^1 \nabla_1 \p_y G(Z_{t_{i-1}} + u \delta_{i-1}^n, Z_{t_{i-1}} + u \delta_{i-1}^n, \theta) \cdot \delta_{i-1}^n \,\df u,
\end{align*}
where $\d_i^n = \int_0^{t_i} [\mu_{\xi_0}(s) - \mu_{\wh{\xi}_n}(s)]\,\df s$, so that the difference is bounded as
\[
\l| \p_y G(Y_{i-1}^n, Y_{i-1}^n, \theta) - \p_y G(Z_{t_{i-1}}, Z_{t_{i-1}}, \theta) \r| \le C(1 + |Z_{t_{i-1}}|^k) |\delta_{i-1}^n|,
\]
for some constant \( C \) and \( k \ge 0 \), using the polynomial growth of \( \nabla_1 \p_y G \).
By Lemma~\ref{lem:ergod}, this yields
\[
\frac{1}{n h_n^2} \sum_{i=1}^n \l| \p_y G(Y_{i-1}^n, Y_{i-1}^n, \theta) - \p_y G(Z_{t_{i-1}}, Z_{t_{i-1}}, \theta) \r| |\Delta_i^n Z| = O_p(h_n) \to 0.
\]
The second term is bounded by
\[
\l| \frac{1}{n h_n^2} \sum_{i=1}^n \p_y G(Y_{i-1}^n, Y_{i-1}^n, \theta) r_i^n \r| \le \frac{C}{n h_n} \sum_{i=1}^n (1 + |Z_{t_{i-1}}|^k) o_p(h_n)= o_p(1),
\]
again using Lemma~\ref{lem:ergod} and the consistency \( \wh{\xi}_n \toP \xi_0 \).
Therefore,  using Lemma \ref{lem:ergod} again, we conclude that 
\[
A_n(\theta) \toP -\frac{\p^2 K(0)}{2K(0)} \int_{\R} \p_y G(z, z, \theta) z \phi_{K(0)}(z)\,\df z, 
\]
The term $B_n(\theta)$ can be evaluated similarly using ergodic theory and standard limit theorems.

As for $C_n(\theta)$, recall that
\[
R_i^n(\theta) = \frac{1}{6} \p_y^3 G\l(X_{t_{i-1}}, \zeta_i, \theta \r) (X_{t_i} - X_{t_{i-1}})^3.
\]
By the assumption that $|\p_y^3 G(x,y,\theta)| \le M$ and the fact that $X_{t_i} - X_{t_{i-1}} = O_p(h_n^{1/2})$, we get
\begin{align*}
|C_n(\theta)| &\le \frac{M}{6 n h_n^2} \sum_{i=1}^n |X_{t_i} - X_{t_{i-1}}|^3 \\
&= O_p\l(\frac{1}{n h_n^2} \sum_{i=1}^n h_n^{3/2} \r) = O_p(h_n^{1/2}) \to 0.
\end{align*}
Hence, the result follows. 
\end{proof}


\begin{remark}
When $G(x,y) = f(y-x)$ for a function $f \in C^2(\mathbb{R})$ with polynomial growth, the limit in Lemma~\ref{lem:GY} simplifies to
\[
\frac{1}{n h_n^2} \sum_{i=1}^n f(Y^n_i - Y^n_{i-1} )
\to^{\mathbb{P}} -\frac{1}{2} \, \partial^2 K(0) \cdot f''(0),
\]
as $n \to \infty$, where $Y^n_i = X_{t_i} - \int_0^{t_i} \mu_{\widehat{\xi}_n}(s)\,\df s$ and $\widehat{\xi}_n \to^{\mathbb{P}} \xi_0$.
This result is consistent with Lemma~\ref{lem:f-DX}, which considers functions of the normalized increment
\[
\frac{\D_i^n X - \mu_\xi(t_{i-1}) h_n}{h_n} \approx \frac{Y^n_{i} - Y^n_{i-1}}{h_n},
\]
and shows convergence of the empirical average of $f$ evaluated at that normalized quantity to the expectation under the normal distribution with mean zero and variance $-\partial_t^2 K(0)$. In this case, the asymptotic mean reduces to a constant multiple of $f''(0)$, matching the second moment of the limiting distribution.
\end{remark}

\begin{remark}
Unlike Lemma~\ref{lem:f-DX}, the convergence in Lemma~\ref{lem:GY} does not generally hold uniformly with respect to $\xi \in \Xi$.
This difference arises from the nature of the function $G(x, y, \theta)$. In Lemma~\ref{lem:f-DX}, the function depends on only the factor $(x - y)$, and the integral term $\int_{t_{i-1}}^{t_i} \mu_\xi(s)\,\df s$ is of order $O(h_n)$ uniformly in $\xi$, allowing the substitution of $\xi$ without affecting the asymptotics. In contrast, the general form of $G(Y^n_{i-1}, Y^n_i, \theta)$ in Lemma~\ref{lem:GY} does not necessarily exhibit such cancellation or stability, and thus the impact of the drift term cannot be neglected. Therefore, it is necessary to incorporate the consistent estimator $\widehat{\xi}_n$ to appropriately correct for the drift component.
\end{remark}






\section{Ergodicity for Gaussian processes}\label{app:C}
\subsection{Fundamental ergodic theorems}
\bi
\item Let $Z$ be a (continuous) Gaussian process on a canonical space $(\Omega,\F,\P)$, where $\Omega = C(\R)$. (So $\P$ is the distribution of $X$)
\item $\th_\tau: \Omega \to \Omega$: the {\it shift operator} such that, for each $\omega\in \Omega$,
\[
\th_\tau\omega(t) = \omega(t - \tau),\quad \tau,t\in \R. 
\]
\ei

\begin{thm}[Maruyama \cite{m49} or Krishnapur \cite{k9}]\label{thm:maruyama}
Let $Z$ be a centered stationary Gaussian process on $\R^d$ with continuous covariance kernel $K$ and the spectral measure $\mu$:
for any $t,s\in \R$,
\[
\E[Z_t] = 0;\quad \E[Z_tZ_s]=K(t-s); \quad K(h) = \int_\R e^{ih x} \mu(\df x).
\]

Then
\bi
\item[(i)] $Z$ is ergodic if and only if $\mu$ has no atom.
\item[(ii) ] $Z$ is weakly mixing if and only if $K(t) =o(1)$ as $|t|\to \infty$.
\ei

\end{thm}

\begin{thm}[Birkoff's ergodic theorem] \label{thm:birkoff}
 $Z$ is ergodic if and only if, for any $g\in L^1(\P)$,
\[
\frac{1}{T}\int_0^T g(\th_\tau \omega)\,\df \tau \to \E[g],\quad T\to \infty, 
\]
almost surely or $L^1(\P)$ sense. Therefore, it converges at least in probability.
\end{thm}

\subsection{Utility format in applications}

\bi
\item Suppose that $Z = (Z_t)_{t\ge 0}\sim GP(0,K)$ with $K(t) \to 0$ as $t\to \infty$.
\item Let $\pi_t\ (t\ge 0)$ is a canonical projection: $\pi_t Z = Z_t$.
\ei

For a function $f:\R\to \R$, of polynomial growth: $|f(x)| \lesssim (1 + |x|)^C$, and a fixed $t\ge 0$,
put $g(\omega) = f\circ \pi_t (\omega)\ (\omega \in \Omega)$, then
\[
g \sim^d f(Z_t) =:f(G),\quad G\sim {\cal N}(0,K(0)), 
\]
by the stationarity, and Thereom \ref{thm:birkoff} says that
\begin{align*}
\frac{1}{T} \int_0^T g(X(\cdot - \tau))\,\df \tau &= \frac{1}{T} \int_0^T f(Z_{T-\tau}) \,\df \tau = \frac{1}{T} \int_0^T f(Z_u) \,\df u \\
&\toP \E[f(G)] = \int_\R f(z) \phi_{K(0)}(z)\,\df z,\quad T\to \infty,
\end{align*}
where $\phi_{\mu, \Sigma}(z)$ is the probability density of ${\cal N}(\mu,\Sigma)$.

\begin{cor}\label{cor:ergodic}
Let $Z = (Z_t)_{t\ge 0}$ be a centered stationary Gaussian process  with $K(t) \to 0\ (t\to \infty)$.
Then it holds for any function $f:\R\to \R$, of polynomial growth that
\begin{align*}
 \frac{1}{T} \int_0^T f(Z_u) \,\df u \to \E[f(Z_0)]\quad \mbox{$as$ or in $L^1$}, \quad T\to \infty,  %\label{Z:ergodic}
\end{align*}
for any measurable function $f:\R\to \R$ such that $\E[f(Z_0)]<\infty$.  
\end{cor}





\begin{thebibliography}{99}

\bibitem{a03} Asmussen, S. (2003). {\it Applied probability and queues}, 2nd ed., Springer-Verlag, New York.

\bibitem{betal24} Bennedsen, M., Christensen, K. and Christensen, P. (2024). Composite likelihood estimation of stationary Gaussian processes with a view toward stochastic volatility. arXiv preprint {\it arXiv:2403.12653v1 [econ.EM]}.

\bibitem{bl20} Bachoc, F. and Lagnoux, A. (2020). Fixed-domain asymptotic properties of maximum composite likelihood estimators for Gaussian processes. {\it J. Statist. Plann. Inference}, {\bf 209}, 62--75.

\bibitem{c12} Caravenna, F. (2012). A note on directly Riemann integrable functions. arXiv preprint {\it arXiv:1210.2361v1 [math.PR]}.

\bibitem{cr04} Cox, D. R. and Reid, N. (2004). A note on pseudolikelihood constructed from marginal densities, {\it Biometrika}, {\bf 91} (3), 729--737.

\bibitem{dy11} Davis, R. A. and Yau, C. Y. (2011). Comments on pairwise likelihood in time series models, {\it Statistica Sinica}, {\bf 21} (1), 255--277.

\bibitem{f71} Feller, W. (1971). {\it An Introduction to Probability Theory and Its Applications}, vol. 2, 2nd ed., John Wiley and Sons, New York.

\bibitem{ft16} Fukasawa, M. and Takabatake, T. (2019). Asymptotically efficient estimators for self-similar stationary Gaussian noises under high-frequency observations. {\it Bernoulli}, {\bf 25} (3), 1870--1900.

\bibitem{gj93} Genon-Catalot, V. and Jacod, J. (1993). On the estimation of the diffusion coefficient for multi-dimensional diffusion processes, {\it Annales de l'I.H.P. Probabilit\'es et statistiques}, {\bf 29} (1), 119--151.


\bibitem{ir78} Ibragimov, I. A. and Rozanov, Yu. A. (1978). {\it Gaussian Random Processes}, Springer-Verlag, New York.

\bibitem{k97} Kessler, M. (1997). Estimation of an ergodic diffusion from discrete observations, {\it Scandinavian Journal of Statistics}, {\bf 24} (2), 211--229.

\bibitem{ketal25} Kobayashi, K., Nishiwaki, Y., Shimizu, Y. and Takaoka, N. (2025). Maximum likelihood estimation of mean functions for Gaussian processes under small noise asymptotics, arXiv preprint {\it arXiv:2507.05628v1 [math.ST]}.

\bibitem{k9} Krishnapur, M. (2014). Lecture Notes for {\it Stationary Gaussian Processes}, {\tt https://math.iisc.ac.in/~manju/GP/gaussianprocesses.html}

\bibitem{m49} Maruyama, G. (1949). The harmonic analysis of stationary stochastic processes, {\it Memoirs of The Faculty of Science, Kyushu University. Series A, Mathematics}, {\bf 4}, 45--106.

\bibitem{n13} Neumann, M. H. (2013). A central limit theorem for triangular arrays of weakly dependent random variables, with applications in statistics. {\it ESAIM: Probability and Statistics}, {\bf 17}, 120--134.

\bibitem{retal99} Rolski, T., Schmidli, H., Schmidt, V. and Teugels, J. (1999). {\it Stochastic processes for insurance and finance}, John Wiley \& Sons, Ltd., Chichester.

\bibitem{rw06} Rasmussen, C. E. and Williams, C. K. I. (2006). {\it Gaussian Processes for Machine Learning}, the MIT Press, Massachusetts Institute of Technology. {\tt www.GaussianProcess.org/gpml}

\bibitem{t22} Takabatake, T. (2024). Quasi-likelihood analysis of fractional Brownian motion with constant drift under high-frequency observations. {\it Statistics \& Probability Letters}, {\bf 207}, 110006.

\bibitem{v98} van der Vaart, A.W. (1998). {\it Asymptotic statistics}. Cambridge University Press, Cambridge.

\bibitem{vetal11} Varin, C., Reid, N. and Firth, D. (2011). An overview of composite likelihood methods, {\it Statistica Sinica}, {\bf 21} (1), 5--42.

\bibitem{w51} Whittle, P. (1953). Estimation and information in stationary time series. {\it Arkiv f\"or Matematik}, {\bf 2} (5), 423--434.



%\bibitem{andas10} Anderes, E. (2010). On the consistent separation of scale and variance for Gaussian random fields. {\it Ann. Statist.}, {\bf 38}, no. 2, 870--893.

%\bibitem{azais20} Aza\"is, J.-M., Bachoc, F. Lagnoux, A. and Nguyen, T. M. N. (2020). Semi-parametric estimation of the variogram scale parameter of a Gaussian process with stationary increments. {\it ESAIM Probab. Stat.}, {\bf 24}, 842--882.

%\bibitem{b14} Bachoc, F. (2014). Asymptotic analysis of the role of spatial sampling for covariance parameter estimation of Gaussian processes. {\it Journal of Multivariate Analysis}, {\bf 125}, 1--35.

%\bibitem{b21} Bachoc, F. (2021). Asymptotic Analysis of Maximum Likelihood Estimation of Covariance Parameters for Gaussian Processes: An Introduction with Proofs. In: {\it Advances in Contemporary Statistics and Econometrics}. Springer, Cham.

%\bibitem{cetal82} Cornfeld, I. P., Fomin, S. V. and Sinai, Y. G. (1982). {\it Ergodic theory}, Springer-Verlag, New York.

%\bibitem{dahlhaus2000} Dahlhaus, R. (2000). A likelihood approximation for locally stationary processes, {\it Annals of Statistics}, {\bf 28} (6), 1762--1794.

%\bibitem{dahlhauspolonik2006} Dahlhaus, R. and Polonik, W. (2006). Nonparametric quasi-maximum likelihood estimation for Gaussian locally stationary processes, {\it Annals of Statistics}, {\bf 34} (6), 2790--2824.

%\bibitem{dd07} Dedecker, J., Doukhan, P., Lang, G., Le\'on, J. R., Louhichi, S. and Prieur, C. (2007). {\it Weak Dependence: With Examples and Applications}. Lecture Notes in Statistics, vol. 190, Springer.

%\bibitem{detal09} Du, J., Zhang, H. and Mandrekar, V. S. (2009). Fixed-domain asymptotic properties of tapered maximum likelihood estimators, {\it The Annals of Statistics}, {\bf 37} (6A), 3330--3361.

%\bibitem{dm02} Dedecker, J. and Merlev\`ede, F. (2002). Necessary and sufficient conditions for the conditional central limit theorem. {\it Annals of Probability}, {\bf 30} (3), 1044--1081.

%\bibitem{dm08} Dym, H. and McKean, H. P. (2008). {\it Gaussian processes, function theory, and the inverse spectral problem}, Dover Publications, New York.

%\bibitem{dmr94} Doukhan, P., Massart, P. and Rio, E. (1994). The functional central limit theorem for strongly mixing processes, {\it Annales de l’Institut Henri Poincar\'e, Probabilit\'es et Statistiques}, {\bf 30} (1), 63--82.

%\bibitem{dp05} Dedecker, J. and Prieur, C. (2005). New dependence coefficients. Examples and applications to statistics. {\it Probab. Theory Relat. Fields}, {\bf 132}, 203--236.

%\bibitem{d95} Doukhan, P. (1995). {\it Mixing: Properties and Examples}. Lecture Notes in Statistics, vol.~85, Springer.

%\bibitem{fgn06} Furrer, R., Genton, M. G. and Nychka, D. (2006). Covariance tapering for interpolation of large spatial datasets. {\it J. Comput. Graph. Statist.}, {\bf 15}, no. 3, 502--523.

%\bibitem{fh23} Furrer, R. and Hediger, M. (2023). Asymptotic analysis of ML-covariance parameter estimators based on covariance approximations. {\it Electron. J. Stat.}, {\bf 17}, no. 2, 3050--3102.

%\bibitem{gs85} Giraitis, L. and Surgailis, D. (1985). CLT and other limit theorems for functionals of Gaussian processes. {\it Zeitschrift f\"ur Wahrscheinlichkeitstheorie und Verwandte Gebiete}, {\bf 70} (2), 191--212.

%\bibitem{il97} Istas, J. and Lang, G. (1997). Quadratic variations and estimation of the local H\"older index of a Gaussian process. {\it Ann. Inst. H. Poincar\'e Probab. Statist.}, {\bf 33}, no. 4, 407--436.

%\bibitem{ko22} Karvonen, T. and Oates, C. J. (2022). Maximum likelihood estimation in Gaussian process regression is ill-posed. arXiv preprint {\it arXiv:2203.09179}.

%\bibitem{l12} Lifshits, M. A. (2012). {\it Lectures on Gaussian processes}, Springer-Verlag, Berlin Heidelberg.

%\bibitem{loh15} Loh, W.-L. (2015). Estimating the smoothness of a Gaussian random field from irregularly spaced data via higher order quadratic variations. {\it Ann. Statist.}, {\bf 43}, no. 6, 2766--2794.

%\bibitem{mm84} Mardia, K.V. and Marshall, R.J. (1984). Maximum likelihood estimation of models for residual covariance in spatial regression, {\it Biometrika}, {\bf 71}, (1), 135--146.

%\bibitem{m17} Minden, V., Damle, A., Ho, K. L. and Ying, L. (2017). Fast spatial Gaussian process maximum likelihood estimation via skeletonization factorizations. {\it Multiscale Modeling \& Simulation}, {\bf 15}(4), 1584--1611.


%\bibitem{pr99} Prakasa Rao, B.L.S. (2010). {\it Statistical Inference for Diffusion Type Processes}. Wiley (The first edition is published by Hodder Headline Group in 1999).


%\bibitem{rllb25} Reding, L., L\'opez-Lopera, A. F. and Bachoc, F. (2025). Covariance parameter estimation of Gaussian processes with approximated functional inputs. {\it J. Multivariate Anal.}, {\bf 205}, 105380.

%\bibitem{s04} Seeger, M. W. (2004). Gaussian processes for machine learning, {\it International journal of neural systems}, {\bf 14} (2), 69--106.

%\bibitem{w94} Wooldridge, J. M. (1994). Estimation and inference for dependent processes, {\it Handbook of Econometrics}: Volume IV, ed. by R. F. Engle, and D. McFadden. North-Holland, Amsterdam, 2639--2738.

%\bibitem{y91} Ying, Z. (1991). Asymptotic properties of a maximum likelihood estimator with data from a Gaussian process, {\it Journal of Multivariate Analysis}, {\bf 36}, 280--296.

\end{thebibliography}
\end{document}
